% =====================================================================
% LES QUESTIONS : partie X et tout ce qui la suit. Separe du cours le 6 octobre 2026.
% =====================================================================
% =====================================================================
% PREAMBULE COMMUN a cours_memoire.tex (parties I a IX) et a
% cours_memoire_questions.tex (partie X et suivantes). Extrait du fichier
% d'origine le 6 octobre 2026. Un seul endroit : toute retouche de style
% vaut pour les deux documents.
% =====================================================================
\documentclass[11pt,a4paper]{article}

\usepackage[utf8]{inputenc}
\usepackage[T1]{fontenc}
\usepackage{lmodern}
\usepackage[french]{babel}
\usepackage[margin=2.1cm]{geometry}
\usepackage{microtype}
\usepackage{amsmath,amssymb}
\usepackage{eurosym}
\usepackage{booktabs}
\usepackage{array}
\usepackage{tabularx}
\usepackage[table]{xcolor}
\usepackage{graphicx}
\usepackage{caption}
\usepackage{enumitem}
\usepackage[most]{tcolorbox}
\usepackage{tikz}
\usetikzlibrary{arrows.meta,positioning,calc,fit,backgrounds,decorations.pathreplacing,shapes.geometric}
\usepackage[hidelinks]{hyperref}

\definecolor{navy}{HTML}{1F3864}
\definecolor{bluemid}{HTML}{2E5496}
\definecolor{lightblue}{HTML}{D6E0F0}
\definecolor{accent}{HTML}{C0491F}
\definecolor{vert}{HTML}{3D8361}

\captionsetup{font=small,labelfont={bf,color=navy},labelsep=period}
\graphicspath{{../figures/}}

\newtcolorbox{phrasebox}{colback=vert!7,colframe=vert!70!black,boxrule=0.7pt,
  arc=2pt,left=7pt,right=7pt,top=5pt,bottom=5pt,before skip=6pt,after skip=8pt}
\newtcolorbox{cle}[1][]{colback=lightblue!40,colframe=navy,boxrule=0.6pt,
  arc=2pt,left=7pt,right=7pt,top=5pt,bottom=5pt,#1}

\providecommand{\Prob}{\mathbb{P}}
\providecommand{\E}{\mathbb{E}}
\providecommand{\R}{\mathbb{R}}

\newcommand{\lequoi}[1]{\smallskip\noindent\textcolor{navy}{\textbf{Le problème.}}\ #1\par}
\newcommand{\lidee}[1]{\smallskip\noindent\textcolor{navy}{\textbf{L'idée.}}\ #1\par}
\newcommand{\laformule}[1]{\smallskip\noindent\textcolor{navy}{\textbf{La formule.}}\ #1\par}
\newcommand{\cheznous}[1]{\smallskip\noindent\textcolor{navy}{\textbf{Chez nous.}}\ #1\par}
\newcommand{\lepiege}[1]{\smallskip\noindent\textcolor{accent}{\textbf{Le piège.}}\ #1\par}
\newcommand{\laphrase}[1]{\begin{phrasebox}\textbf{À dire à voix haute.} \emph{#1}\end{phrasebox}}

% Questions du jury : la question, puis la reponse a tenir.
\newtcolorbox{question}[1]{colback=accent!5,colframe=accent!75!black,boxrule=0.6pt,
  arc=2pt,left=7pt,right=7pt,top=5pt,bottom=5pt,before skip=7pt,after skip=4pt,
  fonttitle=\bfseries\color{white},title={#1}}
\newcommand{\rep}[1]{\smallskip\noindent\textcolor{navy}{\textbf{La réponse.}}\ #1\par}
\newcommand{\appui}[1]{\smallskip\noindent\small\textcolor{navy}{\textbf{Appui~:}} #1\par\normalsize}

% Partie X, chapitre par chapitre : les reperes de pages, le chapitre en une phrase, et la
% fiche \og par coeur\fg{} (tableau a deux colonnes dans un cadre, titre optionnel).
\newcommand{\reperes}[1]{\noindent{\small\color{navy!75}#1}\par\smallskip}
\newcommand{\enbref}[1]{\noindent\textcolor{navy}{\textbf{En une phrase.}}\ #1\par}
\newtcolorbox{parcoeur}[1][À savoir par c\oe ur]{colback=navy!3,colframe=navy!85!black,
  boxrule=0.6pt,arc=2pt,left=6pt,right=6pt,top=4pt,bottom=4pt,before skip=7pt,after skip=5pt,
  fonttitle=\bfseries\small\color{white},title={#1}}
\newenvironment{fichetab}{\small\renewcommand{\arraystretch}{1.1}%
  \tabularx{\linewidth}{@{}>{\raggedright\arraybackslash\bfseries}p{3.2cm}
  >{\raggedright\arraybackslash}X@{}}}{\endtabularx}

% Styles des schemas
\tikzset{
  brique/.style={draw=navy,fill=lightblue!45,rounded corners=2pt,inner sep=5pt,
                 align=center,font=\footnotesize,minimum height=9mm},
  pilier/.style={draw=navy,fill=white,circle,inner sep=1.5pt,minimum size=8mm,
                 font=\footnotesize\bfseries,text=navy},
  fleche/.style={-{Stealth[length=2mm]},draw=navy,line width=0.5pt},
  flechef/.style={-{Stealth[length=2.6mm]},draw=accent,line width=1.1pt},
  flechep/.style={-{Stealth[length=1.8mm]},draw=navy!45,line width=0.4pt},
  etiq/.style={font=\scriptsize\itshape,text=navy!75},
}

\setcounter{secnumdepth}{2}
\emergencystretch=1.2em   % noms propres longs (Kolmogorov-Smirnov, Pickands-Balkema)

\usepackage{longtable}

\title{\vspace{-1.3cm}\color{navy}\bfseries Le rapport ENSAE, chapitre par chapitre\\[3pt]
  \large Ce qu'il faut savoir par c\oe ur, et les questions que le jury peut poser}
\author{Kélian Kaddouri}
\date{Mise à jour du 9 novembre 2026}

\begin{document}
\maketitle
\vspace{-0.8cm}

\begin{cle}
\textbf{Ce document est la suite de \texttt{cours\_memoire\_ensae}.} Ce dernier explique les outils
un par un, des données au capital (parties~I à IX). Celui-ci reprend le rapport ENSAE dans l'ordre
des chapitres, et ne répète ni les schémas ni les démonstrations du cours : si une notion y
paraît obscure, c'est dans le cours qu'il faut la reprendre. \emph{Folios conservés du rapport ENSAE corrigé le 19 octobre ; réponses actualisées au
9 novembre 2026 pour les compléments Institut. Ceux-ci sont regroupés dans une partie finale
hors rapport déposé : les $169$~M\euro{} du rapport ENSAE ne remplacent pas les $24{,}1$~M\euro{}
de la lecture Institut actuelle.}
\end{cle}

\tableofcontents
\newpage

\part*{X. Chapitre par chapitre : ce qu'il faut savoir par c\oe ur}
\addcontentsline{toc}{part}{X. Chapitre par chapitre : ce qu'il faut savoir par c\oe ur}

\begin{cle}
\textbf{Comment lire cette partie.} Une section par chapitre du rapport, dans l'ordre du
document. Chacune donne le chapitre en une phrase, puis un encadré \emph{à savoir par c\oe ur}
(les nombres, définitions et formules à restituer sans notes), puis les questions que le jury
peut poser sur ce chapitre, avec la réponse à tenir et l'endroit du mémoire où s'appuyer. Les
questions qui ne relèvent d'aucun chapitre (forme, déontologie, dispositif) sont regroupées à
la fin.

\textbf{Les folios} sont ceux du rapport ENSAE de $100$~pages (\texttt{memoire/main\_ensae.pdf},
version corrigée du 19 octobre 2026) : introduction folio~$11$, chapitres~$1$ à~$10$ des folios~$13$ à~$77$,
conclusion folio~$78$, enseignements du stage folio~$80$, annexe~A folio~$87$, annexe~B folio~$91$,
annexe~C folio~$95$, notes de synthèse folio~$96$. Le \emph{supplément} (S1 à S4, DOI
\texttt{10.5281/zenodo.22874524}) est hors du rapport : le jury de l'ENSAE ne l'a pas sous les yeux.
La numérotation des chapitres est celle du rapport ENSAE : les anciens chapitres~$4$ à~$11$ du
manuscrit complet y sont les chapitres~$3$ à~$10$, le cadre prudentiel est la section~$1.6$ et la
donnée manquante la section~$10.8$.

\textbf{Trois règles pour l'oral ENSAE} : tenir \textbf{15 minutes au maximum}, questions séparées ;
expliquer la méthode et la portée des résultats ; \emph{si vous ne savez pas, dites-le}.
La durée Institut doit être confirmée dans sa convocation : une cible de préparation ne vaut pas consigne officielle.
\end{cle}

\begin{cle}
\textbf{Quatre réflexes qui évitent les fautes graves.} (i)~Dire \og au secteur\fg{} avant tout
montant en milliers de millions. (ii)~Dire \og besoin de capital ORSA au titre de DORA\fg{},
jamais \og le SCR DORA de telle entité\fg{}. (iii)~Ne jamais additionner les quatre canaux, ni
l'interaction entre canaux avec celle entre piliers. (iv)~Ne jamais comparer un quantile de
sévérité d'un sinistre ($663$~M\euro) à un capital ($6\,049$ ou $8\,123$~M\euro) : ce ne sont
pas des objets de même échelle.
\end{cle}

% ---------------------------------------------------------------------
\section*{Les pages liminaires : résumé, abstract et notes de synthèse}
\addcontentsline{toc}{section}{Les pages liminaires : résumé, abstract et notes de synthèse}
\reperes{Résumé folio~3, abstract folio~4, glossaire folio~9, note de synthèse folio~96, executive
summary folio~98.}

\enbref{Ces pages sont lues en premier, et parfois seules : chaque nombre qu'elles portent doit
sortir sans hésitation, et dans la même formulation que le résumé.}

\begin{parcoeur}
\begin{fichetab}
UE 2022/2554 & le règlement DORA, applicable depuis le 17 janvier 2025 \\
$1\,024$ sommets & les $2^{10}$ orientations extrêmes des dix paires de piliers \\
$[6\,858\,;8\,697]$ M\euro & plage aux sommets, conditionnelle à $S$ posée ;
  sa largeur de $1\,839$ mesure l'effet des directions évaluées \\
$6\,049\to20\,188$ M\euro & conforme vers non conforme, quatre canaux, au secteur :
  facteur $3{,}34$, écart $14\,139$ \\
$9\,138$ et $5\,001$ & la somme des canaux isolés, puis leur interaction ($35\,\%$ de
  l'écart) \\
$10\,377$ et $11\,413$ & la fréquence seule, puis les trois autres canaux réunis \\
$131{,}5$ à $183{,}3$ M\euro & l'entité notionnelle, extrapolation d'ordre de grandeur ($169$ au
  point) \\
facteur $22$, $41$ & la variation du niveau avec la seule source de sévérité, puis avec la
  fréquence triplée en plus \\
$32{,}7\,\%$ et $13{,}3\,\%$ & la réduction moyenne de la plage si le sens de six dépendances était
  établi ($22$ à $44\,\%$ selon la graine), puis par la seule paire gouvernance-incidents (quatre
  paires entre $11{,}2$ et $13{,}3\,\%$, ordre non résolu) \\
$2{,}78$ & le coût annuel de détention rapporté au prix de marché du transfert (taux révisé) \\
$1\,811$ années & l'historique nécessaire pour backtester le quantile à $99{,}5\,\%$ \\
\end{fichetab}
\end{parcoeur}

\begin{question}{\og Résumez votre mémoire en une minute.\fg{}}
\rep{Je calcule ce que la non-conformité à DORA coûte en capital. Je modélise les cinq piliers avec une cascade : quand un pilier tombe, il peut entraîner les autres. Un point reste inconnu : le sens de la cascade. Les données publiques ne le donnent pas. Je borne donc le capital, sur les $1\,024$ configurations admissibles, entre $6\,858$ et $8\,697$~millions. L'écart entre ces deux valeurs est le prix de cette ignorance. Ensuite, la conformité change quatre choses dans le modèle. Le besoin de capital du secteur passe alors de $6\,049$ à $20\,188$~millions, soit $3{,}34$ fois plus. C'est surtout la fréquence des incidents qui joue. Enfin, je dis quelle donnée il faudrait réclamer et ce qu'elle apporterait.}
\appui{résumé, folio 3.}
\end{question}

\begin{question}{\og La note parle d'un facteur $22$ sur le niveau, le résumé d'un facteur $41$. Lequel ?\fg{}}
\rep{Les deux sont justes : ce sont deux variations différentes. Le facteur $22$ vient d'un seul changement, celui de la source de sévérité (OpRisk ou PRC). Le facteur $41$ vient du même changement, avec en plus une fréquence multipliée par trois. Dans les deux cas, le message est le même : le niveau dépend de choix qu'on ne peut pas transposer d'une base à l'autre. En revanche, les rapports au socle varient de $33\,\%$ au plus.}
\appui{\S~8.4, folio 58 ; note de synthèse, rubrique Limites.}
\end{question}

\begin{question}{\og Deux canaux reposent sur des bornes posées, écrivez-vous. Lesquels, et qu'est-ce que cela change ?\fg{}}
\rep{Les deux canaux posés sont la propagation ($g$ de $0{,}45$ à $0{,}90$) et l'accumulation par un prestataire commun (de $0$ à $0{,}68$). Pour la fréquence et la détection, le niveau de départ est calibré sur les données. Mais l'écart entre les deux états est posé pour les quatre canaux. Pour la fréquence, il vient d'effets publiés. Pour les deux canaux posés, l'effet de fermeture est une borne. Ce n'est pas un gain de remédiation garanti. Leur part de Shapley n'est pas non plus un budget.}
\appui{\S~6.3, folios 46 et 47 ; supplément S1, \S~S1.4.}
\end{question}

% ---------------------------------------------------------------------
\section*{Chapitre 1. Risque cyber, DORA et cadre prudentiel (sections 1.1 à 1.5)}
\addcontentsline{toc}{section}{Chapitre 1. Risque cyber, DORA et cadre prudentiel (sections 1.1 à 1.5)}
\reperes{folios 13 à 19. 1.1 le risque (13), 1.2 l'assurabilité (14), 1.3 le marché français (15), 1.4 DORA (15), 1.5 l'enjeu prudentiel et la problématique (17).}

\enbref{Le cyber n'est pas un risque de plus : sa charge est portée par quelques événements, ses
causes se propagent par les prestataires partagés, et le marché français n'a pas l'historique
pour en estimer la queue. Ces trois traits commandent la théorie des valeurs extrêmes, la cascade
dirigée et la calibration sur une base internationale.}

\begin{parcoeur}
\begin{fichetab}
Incident, attaque & l'incident TIC est tout événement qui affecte la sécurité ; l'attaque exige
  trois conditions cumulatives (intention, impact significatif, aboutissement). DORA retient
  l'incident : la panne CrowdStrike de 2024 est dans son champ \\
Quatre familles & rançongiciel, fuite de données, déni de service, défaillance d'un prestataire,
  la seule non individualisable, qui produit des sinistres corrélés \\
Quatre cas & WannaCry 2017 (échelle), NotPetya 2017 (P4 puis P1 puis P2), MOVEit 2023
  (accumulation), CrowdStrike 2024 (dépendance sans adversaire, P4 et P3) \\
LUCY 2026 & $20\,996$ polices, $1\,251$ sinistres, sept exercices ; assurés $+49\,\%$, primes
  $-3\,\%$ ; sinistres sur primes $12\,\%$ en 2023, $27\,\%$ en 2025, $169\,\%$ en 2020 ; charge
  indemnisée $83{,}2$ M\euro{} \\
Deux régimes & 2024, sévérité (nombre $\times0{,}73$, coût moyen $\times1{,}96$) ; 2025,
  fréquence (nombre $\times2{,}79$, coût moyen $\times0{,}55$) ; un seul sinistre au-delà de
  $10$ M\euro{} en 2025 \\
Cinq piliers & P1 gouvernance (art.~5 à 16), P2 incidents (17 à 23), P3 tests (24 à 27, TLPT
  de l'art.~26 tous les trois ans), P4 tiers (28 à 44, clauses de l'art.~30, troisième rang de
  sous-traitance), P5 partage (art.~45) \\
Notification & alerte sous 4~h après classification et au plus 24~h après détection, rapport
  intermédiaire sous 72~h, rapport final sous un mois \\
AMF 2025 & $47$ notifications, $31$ incidents majeurs, $27$ d'origine tierce ($87\,\%$) : une
  part d'origine, pas une probabilité de transfert \\
CESIN, France Ass. & attaque par un tiers $35\,\%$, $43\,\%$ dans les grandes entreprises ;
  sanction d'une autorité dans $3\,\%$ des cas ; cyber premier risque pour la huitième année
  ($4{,}1$ sur $5$), fréquence $-0{,}45$ et sévérité $+0{,}25$ en 2026 \\
\end{fichetab}
\end{parcoeur}

\begin{question}{\og Pourquoi DORA, et pas NIS~2, le RGPD ou le règlement sur la cyberrésilience ?\fg{}}
\rep{La mission porte sur la résilience numérique des entités financières. Les cinq domaines de DORA fournissent la grille de contrôles du modèle. Les autres textes précisent le contexte, mais ne sont pas l'objet de la quantification menée ici. Cette grille ne rend pas automatiquement les défaillances observables : leur codage et les données disponibles doivent encore être examinés.}
\appui{\S~1.4, folios 15 à 17.}
\end{question}

\begin{question}{\og Vous êtes co-auteur de l'étude LUCY. N'est-ce pas gênant de la citer ?\fg{}}
\rep{Je la cite comme une source externe que je ne peux pas recalculer. Aucune de ses valeurs n'entre dans un calcul. Aucune ne sert à fixer un capital. Ce qui peut être vérifié l'est : un script contrôle que mes recopies sont fidèles. Il vérifie les sinistres sur primes, la somme des classes de taille et les multiplicateurs. Enfin, trois de ses constats servent le mémoire sans le flatter. Il n'existe pas de traduction prudentielle. La queue nationale est vide. Et la charge se décompose entre nombre de sinistres et coût moyen.}
\appui{\S~1.3, folio 15.}
\end{question}

\begin{question}{\og Vos milliers de millions contre les $83$~millions indemnisés du marché français : lequel est faux ?\fg{}}
\rep{Aucun des deux n'est faux : ils ne mesurent pas la même chose. LUCY publie une charge \emph{indemnisée}. Elle est nette de franchise, plafonnée par la capacité, et elle porte sur le marché français. Moi, je modélise la perte \emph{brute} d'entités financières, à l'échelle du secteur mondial. Si on les compare, on prend une différence de définition pour une contradiction.}
\appui{\S~1.3, folio 15, encadré.}
\end{question}

\begin{question}{\og Le marché voit presque tripler le nombre de sinistres, les enquêtes voient la fréquence baisser. Qui croire ?\fg{}}
\rep{Les deux sont justes. Le nombre de sinistres du marché est un total. Il augmente parce que la base d'assurés a grandi de $49\,\%$. Le CESIN et la cartographie mesurent, eux, l'exposition d'une entité. On peut compter plus de sinistres tout en assurant des entités moins exposées chacune. C'est pourquoi je parle de \emph{nombre} de sinistres, et non de fréquence, pour les multiplicateurs de LUCY.}
\appui{\S~1.5, folios 17 et 18, encadré.}
\end{question}

\begin{question}{\og Le vrai coût de la non-conformité, n'est-ce pas l'amende ?\fg{}}
\rep{Non, et les données le montrent. Une sanction d'une autorité n'apparaît que dans $3\,\%$ des cas du CESIN. C'est loin derrière l'arrêt de la production ($28\,\%$) et la perte d'image ($26\,\%$). De plus, mon modèle n'ajoute aucune pénalité. La non-conformité change seulement quatre paramètres de la loi de perte. Faire de l'amende le coût principal, ce serait se tromper de mécanisme.}
\appui{\S~1.5, folios 17 et 18 ; \S~6.3, folios 46 et 47.}
\end{question}

\begin{question}{\og L'AMF trouve $87\,\%$ d'incidents d'origine tierce. Pourquoi ne pas calibrer P4 dessus ?\fg{}}
\rep{Parce que c'est une part d'\emph{origine}, et non une probabilité de transfert. Elle ne dit pas comment un domaine de contrôle en entraîne un autre. Elle confirme une chose : un prestataire compromis touche plusieurs sociétés à la fois. Elle ne dit rien de la force du lien entre P4 et P2. Enfin, l'écart entre $47$ notifications et $31$ incidents confirmés montre ce que sera le registre. Ce sera une collecte au-dessus d'un seuil de gravité, pas un recensement complet.}
\appui{\S~1.4, folios 15 à 17.}
\end{question}

\begin{question}{\og Pourquoi la Formule Standard ne suffit-elle pas ?\fg{}}
\rep{La Formule Standard charge le risque opérationnel par un pourcentage des primes ou des provisions. Deux entités de même bilan, l'une très mature sur les cinq piliers et l'autre pas du tout, paient donc la même charge. Il n'existe aucun sous-module cyber. Une loi jointe des seules pertes ne décrit pas, à elle seule, l'ordre des transmissions. Le comparateur à copule gaussienne figée testé ici ne fait pas varier cette dépendance avec l'état ; cette limite porte sur sa spécification, pas sur toute copule.}
\appui{\S~1.5, folios 17 et 18.}
\end{question}

\lepiege{Les tentatives WannaCry bloquées (plus de $64$~millions en Russie) ne sont pas des
sinistres : elles mesurent la persistance d'une exposition, et ne se comparent à aucune grandeur
du modèle.}

% ---------------------------------------------------------------------
\section*{Section 1.6. Le cadre prudentiel de capital et son articulation avec DORA}
\addcontentsline{toc}{section}{Section 1.6. Le cadre prudentiel de capital et son articulation avec DORA}
\reperes{folios 18 et 19, section 1.6.}

\enbref{Le risque opérationnel a la même définition en banque et en assurance mais pas le même
traitement : Bâle a retiré les modèles internes, Solvabilité~II garde la liberté de méthode, et
c'est dans l'ORSA, non dans le forfait du pilier~1, que se loge le mémoire.}

\begin{parcoeur}
\begin{fichetab}
Risque opérationnel & perte due à des procédures, des personnes ou des systèmes inadaptés ou
  défaillants, ou à des événements extérieurs ; inclut le juridique, exclut le stratégique et la
  réputation \\
Formule Standard & $\mathrm{SCR}=\mathrm{BSCR}+\mathrm{Adj}+\mathrm{SCR}_{Op}$, avec
  $\mathrm{SCR}_{Op}=\min(0{,}3\,\mathrm{BSCR},\,\mathrm{Op})+0{,}25\,\mathrm{Exp}_{ul}$ et
  $\mathrm{Op}=\max(\mathrm{Op}_{\text{primes}},\mathrm{Op}_{\text{provisions}})$ ; article~204
  du règlement délégué 2015/35 ; terme de croissance au-delà de $20\,\%$ de hausse des primes \\
Bâle & modèles internes AMA retirés en 2022 au profit de l'approche standardisée ; l'approche
  par distribution des pertes reste un outil de modélisation \\
Solvabilité~II & Formule Standard, paramètres propres (USP), ou modèle interne partiel ou
  global approuvé \\
Piliers 1 et 2 & le forfait du pilier~1, plafonné à $30\,\%$ du BSCR, contre l'ORSA du
  pilier~2, prospectif, qui admet les scénarios de stress stochastiques \\
Additivité & $\mathrm{SCR}_{Op}$ n'entre pas dans la matrice de corrélation : un euro d'écart
  DORA vaut un euro de capital (\S~10.12) \\
\end{fichetab}
\end{parcoeur}

\begin{question}{\og Écrivez la charge de risque opérationnel de la Formule Standard.\fg{}}
\rep{$\mathrm{SCR}=\mathrm{BSCR}+\mathrm{Adj}+\mathrm{SCR}_{Op}$, avec
$\mathrm{SCR}_{Op}=\min(0{,}3\,\mathrm{BSCR},\mathrm{Op})+0{,}25\,\mathrm{Exp}_{ul}$. Ici, $\mathrm{Op}$ est le plus grand de deux montants : un pourcentage des primes ou un pourcentage des provisions. Cette charge ne dépend que du volume. Elle ne réagit donc pas à la conformité, par construction.}
\appui{\S~1.6, folios 18 et 19.}
\end{question}

\begin{question}{\og Pourquoi l'ORSA, et pas un modèle interne de pilier~1 ?\fg{}}
\rep{Aucun module DORA n'existe dans la Formule Standard. Un modèle interne exigerait une validation et une donnée d'entité que personne n'a. L'ORSA, lui, demande d'évaluer si le profil de risque réel est couvert par les fonds propres. Il accepte aussi des scénarios stochastiques. Ce que je calcule est donc un besoin de capital du pilier~2. Ce n'est pas un SCR.}
\appui{\S~1.6, folios 18 et 19.}
\end{question}

\begin{question}{\og Bâle a abandonné les modèles internes. Pourquoi y revenir ?\fg{}}
\rep{Bâle a retiré les modèles internes du calcul réglementaire. Mais l'approche par distribution des pertes reste un bon outil de modélisation. J'en reprends les deux marges, la fréquence et la sévérité, sur la base qui servait à ces modèles. Je remplace seulement leur façon d'agréger par la cascade. Et l'assurance garde, sous Solvabilité~II, une liberté de méthode que la banque a perdue.}
\appui{\S~1.6, folios 18 et 19 ; \S~3.7, folios 27 et 28.}
\end{question}

\begin{question}{\og Pourquoi un euro de DORA vaudrait-il plus qu'un euro de risque diversifiable ?\fg{}}
\rep{Parce que la charge opérationnelle s'ajoute au BSCR sans passer par la matrice de corrélation. Un écart est donc reconnu en entier. Si on l'agrégeait comme un module corrélé à $\rho$, il ne compterait que pour $\tau\approx\rho$. À $\rho=0{,}25$, l'écart de traitement va de $3{,}4$ à $3{,}9$ sur les quatre entités réelles. Cette conclusion ne dépend d'aucun paramètre de la cascade. Elle dépend seulement de la place que le règlement donne au risque opérationnel.}
\appui{\S~9.12, folios 69 et 70.}
\end{question}

\lepiege{Ne pas dire que la Formule Standard sous-estime le cyber : elle ne le voit pas. Et le
repère de $3\,\%$ des provisions employé au chapitre~10 n'est qu'un ordre de grandeur, qui
s'écarte de $-8\,\%$ et $+56\,\%$ des modules publiés.}

% ---------------------------------------------------------------------
\section*{Chapitre 2. Revue de littérature et positionnement}
\addcontentsline{toc}{section}{Chapitre 2. Revue de littérature et positionnement}
\reperes{folios 20 à 23. 2.1 la signature (20), 2.2 le positionnement (21), 2.3 la revue par thème (21), 2.4 où se situe le mémoire (23).}

\enbref{Trois vagues de littérature (assurabilité et valeurs extrêmes, contagion entre entités,
Hawkes dans le temps) ; le mémoire en garde les marges et ajoute un mécanisme qu'aucune ne porte,
une dépendance dirigée et sensible à l'ordre entre les domaines de contrôle d'une entité,
traduite en capital.}

\begin{parcoeur}
\begin{fichetab}
Signature & sur les $26$ sous-ensembles d'au moins deux piliers, $22$ changent de criticité
  selon l'ordre ; écart jusqu'à $6$ points sur $10$ ; la chaîne P1 vers P2 vaut $8$, P2 vers P1
  vaut $2$ ($T_{12}=0{,}80$, $T_{21}=0{,}20$) \\
Les concurrents & AMDEC (ni direction ni propagation), copule figée (ni transmission ni ordre), réseau et
  épidémiologie (entre entités, exige la structure du réseau), Hawkes (dans le temps, entre types
  d'attaque, exige un horodatage fin), réseau bayésien (acyclique, tables conditionnelles au même
  mur d'identification) \\
Emprunts & au crédit, le branchement multitype et les modèles multi-états de migration ; à
  Leontief, la normalisation ; à Hillairet et Lopez, les états de conformité ; à Manski (1990,
  1995) et Tamer (2010), l'identification partielle \\
Précédents & Eling et Wirfs (2019), sur la même base SAS OpRisk, pour la sévérité ; Ranguin
  (2024), mémoire pour la CCR sur des sources proches, sans capital à un quantile réglementaire \\
Préprint climatique & entièrement synthétique : on lui emprunte une architecture, aucun chiffre \\
Non revendiqué & la non-transitivité, testée dans trois formulations, aucune ne tient \\
\end{fichetab}
\end{parcoeur}

\begin{question}{\og Qu'est-ce qui est vraiment nouveau dans votre travail ?\fg{}}
\rep{La nouveauté, c'est le mécanisme de dépendance. Il est dirigé et il dépend de l'ordre entre les piliers réglementaires d'une même entité. Je le traduis en capital. Je démontre sa frontière d'identification, puis j'en tire une spécification de reporting. Je n'invente ni loi de sévérité ni processus de fréquence. Pour les marges, je calibre ce que dit la littérature. Je suis donc ambitieux sur le mécanisme, sobre sur le reste, et clair sur les limites.}
\appui{\S~2.4, folio 23, encadré.}
\end{question}

\begin{question}{\og Pourquoi pas un réseau bayésien, qui porte aussi une direction ?\fg{}}
\rep{Un réseau bayésien suppose un graphe sans cycle. Or la contagion entre piliers a des cycles, et la forme de Leontief sait les gérer. De plus, ses tables de probabilités conditionnelles butent sur le même mur que $W$ : on ne peut pas identifier la direction. Le réseau bayésien déplacerait le problème sans le résoudre.}
\appui{\S~2.2, folio 21.}
\end{question}

\begin{question}{\og Quel est le travail le plus proche du vôtre ?\fg{}}
\rep{Pour la sévérité, c'est Eling et Wirfs (2019). Ils étudient les pertes cyber par les valeurs extrêmes, sur la même base SAS OpRisk. Pour les sources, c'est le mémoire de Ranguin (2024). Il porte sur la sinistralité d'un portefeuille assuré. Il ne mesure pas de capital à un quantile réglementaire. Il ne traite pas non plus la dépendance entre domaines de contrôle.}
\appui{\S~2.3, folios 21 à 23.}
\end{question}

\begin{question}{\og Qu'empruntez-vous aux travaux de Caroline Hillairet ?\fg{}}
\rep{J'emprunte d'abord la logique d'états de conformité de Hillairet et Lopez. Je l'applique entre les piliers d'une entité, et non entre des assurés. J'emprunte aussi des noyaux pour tester le processus de Hawkes : ceux de Bessy-Roland, Boumezoued et Hillairet, puis de Boumezoued, Cherkaoui et Hillairet. Je le teste donc contre leurs vrais noyaux, pas contre un noyau facile à battre.}
\appui{\S~2.2 et 2.3, folios 21 à 23 ; \S~10.3, folios 73 et 74.}
\end{question}

\begin{question}{\og Ce préprint climatique qui confirme vos résultats, que vaut-il ?\fg{}}
\rep{Ce préprint est synthétique, et il le dit lui-même. Aucun de ses chiffres n'est un repère empirique, et je n'en reprends aucun. J'emprunte seulement une architecture : la séparation des échelles de temps, les marges appariées et l'ablation. Nos classements diffèrent à un endroit. La cause est dans son texte : sa sévérité est bornée, donc il n'a pas de brique de queue. Chez moi, la queue est la brique la plus lourde.}
\appui{\S~2.3, folios 21 à 23, encadré ; \S~9.7, folios 65 et 66.}
\end{question}

% ---------------------------------------------------------------------
\section*{Chapitre 3. Données}
\addcontentsline{toc}{section}{Chapitre 3. Données}
\reperes{folios 24 à 30. 3.1 cartographie (24), 3.2 trois biais (24), 3.3 actualisation (25), 3.4 trois fréquences (26), 3.5 deux biais chiffrés (26), 3.6 Hackmageddon (27), 3.7 OpRisk (27), 3.8 profilage des variables (28), 3.9 portée (30), 3.10 pourquoi pas plus de sources (30).}

\enbref{Aucune base ne dit tout : chaque source porte ce qu'elle peut porter, ses biais sont
déclarés puis mesurés quand c'est possible, et la queue commune est une décision adossée à des
tailles d'effet, pas une commodité.}

\begin{parcoeur}[À savoir par c\oe ur : les sources et les fréquences]
\begin{fichetab}
PRC & $74\,783$ incidents dans la chronologie, $15\,053$ retenus sur 2019--2025 ; sévérité
  \emph{dérivée} par Jacobs, $\ln L=7{,}68+0{,}76\ln X$ ; $\hat\xi_{\text{PRC}}=1{,}03$
  $[0{,}97\,;1{,}09]$ \\
SAS OpRisk & $41\,975$ pertes ; $582$ cyber $\times$ finance sur 2000--2026 pour la GPD ;
  $1\,476$ TIC en finance sur 2005--2022 pour la fréquence et la descente ; le décile supérieur
  porte $92{,}3\,\%$ de la perte, le centile $65{,}1\,\%$ \\
Hackmageddon & $1\,041$ incidents de janvier à juin 2026, $840$ à vecteur identifié ($81\,\%$) ;
  hameçonnage $38{,}8\,\%$, exploitation $33{,}8\,\%$, chaîne d'approvisionnement $15{,}8\,\%$,
  identifiants $6{,}3\,\%$ ; citation, non recalculable \\
Trois fréquences & $\lambda_{ref}=341$/an (notifications PRC des organisations financières,
  code BSF ; $2\,150$ sur la base entière) ; $\lambda=21{,}6$/an (incidents TIC matériels du
  secteur, seule entrée du moteur) ; $\lambda_{\text{entité}}=0{,}092$/an (lu à la taille de
  l'entité ; les seaux valent $0{,}10$ et $0{,}21$) ; facteur $3\,719$ entre les extrêmes \\
Montants & nominaux ; inflation implicite $2{,}19\,\%$ par an (rapport $1{,}747$ en 2004), non
  appliquée à la calibration gelée \\
\end{fichetab}
\end{parcoeur}

\begin{parcoeur}[À savoir par c\oe ur : les biais et le profilage]
\begin{fichetab}
Trois biais déclarés & sélection (OpRisk, grandes entités), déclaration (PRC, les fuites plus
  que les rançongiciels sans exfiltration), survie (la ruine sort de la base : queue allégée) \\
Deux biais mesurés & élasticité sévérité-taille corrigée de la troncature $b=0{,}087$
  $[0{,}026\,;0{,}148]$, soit une sévérité multipliée par $0{,}85$ à l'entité ; chain-ladder :
  $97$ événements sur 2023--2025 deviennent $255$ à l'ultime \\
Compensation & $0{,}861\times1{,}147=0{,}987$ sur la perte moyenne ; pas sur le quantile \\
Complétude & un manquant PRC vaut la chaîne \texttt{UNKN} : un contrôle naïf donne $100\,\%$ ;
  le volume touché n'est renseigné qu'à $41{,}77\,\%$, et les lacunes se cumulent \\
Doublons & $1\,643$ lignes, $10{,}9\,\%$, borne basse ; le mémoire n'additionne jamais ces
  volumes \\
Queue commune & $V$ de Cramér $<0{,}20$ (max $0{,}176$), $\varepsilon^2<0{,}08$ (max
  $0{,}0785$), Spearman $<0{,}16$ (actif total $0{,}068$, dernier des trois) \\
Seuil & $20{,}03$ M\euro{}, percentile $84{,}4$ de la donnée courante ($q_{85}=22{,}03$) ;
  exactement $91$ excès ; maximum $=3{,}7\times q_{99}$ contre $52{,}7$ sur la base entière \\
\end{fichetab}
\end{parcoeur}

\begin{question}{\og Pourquoi calibrer une entité française sur des pertes de grandes banques mondiales ?\fg{}}
\rep{Parce que la base française est vide là où se joue le quantile. LUCY ne compte qu'un seul sinistre au-delà de $10$~M\euro{} en 2025. Sa classe la plus haute est vide deux années sur sept. Le prix à payer est un biais de taille. Je le mesure au lieu de le supposer : $b=0{,}087$. La sévérité transposée à l'entité n'est multipliée que par $0{,}85$. C'est une réserve d'extrapolation, pas une preuve que le chiffre d'entité majore le besoin réel.}
\appui{\S~3.5, folios 26 et 27 ; \S~1.3, folio 15.}
\end{question}

\begin{question}{\og Votre fréquence vient d'une base américaine sans seuil de matérialité.\fg{}}
\rep{Oui, c'est vrai pour $\lambda_{ref}=341$. C'est pour cela qu'il ne fixe aucun niveau. Il sert seulement à répartir la fréquence entre les vecteurs d'attaque. Le moteur utilise $\lambda=21{,}6$. Ce sont les incidents TIC \emph{matériels} lus dans OpRisk, au-dessus d'un seuil de collecte.}
\appui{\S~3.4, folio 26, encadré ; \S~4.3, folios 33 et 34.}
\end{question}

\begin{question}{\og Trois nombres s'appellent fréquence et vont de $0{,}092$ à $341$. Lequel est le bon ?\fg{}}
\rep{Les trois sont justes. Ce sont trois niveaux d'observation : des notifications, des pertes matérielles du secteur, et les mêmes pertes ramenées à une firme de la taille visée. Le facteur $3\,719$ entre les deux extrêmes est recalculé à chaque exécution. Un seul des trois alimente le moteur : $\lambda=21{,}6$.}
\appui{\S~3.4, folio 26, encadré.}
\end{question}

\begin{question}{\og Comment avez-vous traité les biais de vos bases ?\fg{}}
\rep{Trois biais sont déclarés avant usage : la sélection, la déclaration et la survie. Deux sont mesurés. Leurs corrections vont en sens opposés : $0{,}861\times1{,}147=0{,}987$. L'effet net est de $1{,}3\,\%$ sur la perte moyenne. Mais elles ne se compensent pas sur le capital. Avec une queue lourde, le quantile suit la sévérité. Or la correction de taille porte justement sur la sévérité.}
\appui{\S~3.2 et 3.5, folios 24 à 27.}
\end{question}

\begin{question}{\og Pourquoi une seule queue pour toutes les catégories ?\fg{}}
\rep{Parce que trois outils indépendants disent la même chose : la catégorie informe peu sur la sévérité. Le $V$ de Cramér est sous $0{,}20$. Le $\varepsilon^2$ est sous $0{,}08$. Les corrélations de rang sont sous $0{,}16$. De plus, couper par les sept catégories bâloises laisserait environ douze excès par catégorie parmi les $91$. On ne peut pas ajuster une loi de queue sur si peu. Choisir une queue commune est donc une mesure, pas une commodité.}
\appui{\S~3.8, folios 28 et 29 ; supplément S1, \S~S1.1.}
\end{question}

\begin{question}{\og Vos montants ne sont pas actualisés.\fg{}}
\rep{C'est exact. La chaîne lit des montants nominaux. L'inflation implicite de la base est de $2{,}19\,\%$ par an. Je ne l'applique pas à la calibration gelée, ce qui sous-estime la queue. Je mesure son effet au backtest : la moitié de la dérive de $4{,}00\,\%$ par an est de l'inflation.}
\appui{\S~3.3, folio 25 ; \S~4.7, folio 36.}
\end{question}

\begin{question}{\og Une de vos sources n'est pas reproductible.\fg{}}
\rep{Oui, c'est Hackmageddon. Je la déclare comme une citation, pas comme une sortie de mon dispositif. C'est aussi la source dont le modèle dépend le moins. Elle donne les parts par vecteur d'attaque, donc une structure, jamais un niveau. La catégorie d'un incident ne sert pas à segmenter. Son statut de citation s'étend aussi au multiplicateur de fréquence, et je le déclare.}
\appui{\S~3.6, folio 27, encadré.}
\end{question}

\begin{question}{\og Votre queue n'est-elle pas tirée par un sinistre hors norme ?\fg{}}
\rep{Non. Sur le périmètre calibré, le plus gros sinistre ne vaut que $3{,}7$ fois le centile $99$. Sur la base entière, il en vaut $52{,}7$. Le plus gros sinistre du fichier n'est pas un risque informatique. Et la plus grosse perte calibrée fait partie d'une suite continue de pertes, sans saut isolé.}
\appui{\S~3.8, folios 28 et 29, encadré ; \S~4.6, folios 35 et 36.}
\end{question}

\begin{question}{\og Pourquoi ne pas ajouter d'autres bases ?\fg{}}
\rep{Chaque base apporterait ses propres biais. Elle ne lèverait pas la limite de fond : il n'existe aucun registre harmonisé. De plus, c'est l'indice de queue qui domine l'incertitude. Mieux exploiter les sources existantes, par bootstrap et validation, réduit plus l'incertitude qu'une quatrième base.}
\appui{\S~3.10, folio 30.}
\end{question}

\lepiege{Ne jamais présenter $\hat\xi_{\text{PRC}}=1{,}03$ comme un risque plus lourd : la
sévérité PRC est une puissance du nombre d'enregistrements, et une puissance multiplie l'indice
de queue par son exposant. Les deux indices ne portent pas sur le même objet.}

% ---------------------------------------------------------------------
\section*{Chapitre 4. Modèle fréquence-sévérité}
\addcontentsline{toc}{section}{Chapitre 4. Modèle fréquence-sévérité}
\reperes{folios 31 à 36. 4.1 les deux théorèmes (31), 4.2 estimer $\xi$ (32), 4.3 la fréquence (33), 4.4 la sévérité (34), 4.5 Hill contre maximum de vraisemblance (34), 4.6 validation (35), 4.7 backtest (36).}

\enbref{Deux lois reprises de la littérature et calibrées : une binomiale négative pour le nombre
d'incidents, une loi de Pareto généralisée au-delà de $20{,}03$~M\euro{} pour la sévérité. La
validation établit les formes, pas le niveau, et le backtest dit en années ce qu'il ne peut pas
tester.}

\begin{parcoeur}[À savoir par c\oe ur : la sévérité]
\begin{fichetab}
Deux théorèmes & Fisher-Tippett-Gnedenko : le maximum renormalisé tend vers une GEV ;
  Pickands-Balkema-de Haan : les excès au-delà d'un seuil tendent vers une GPD de même $\xi$ \\
Moments & $\xi>0$ queue lourde ; variance infinie dès $\xi\ge1/2$, espérance infinie dès
  $\xi\ge1$ \\
Calibration & $u=20{,}03$ M\euro{}, $91$ excès, $\hat\xi=0{,}5954$, $\hat\sigma=57{,}97$,
  $\zeta_u=0{,}151$ (gelé) \\
Quantile & $q_\alpha=u+\frac{\sigma}{\xi}\bigl[\bigl(\frac{1-\alpha}{\zeta_u}\bigr)^{-\xi}-1\bigr]$,
  soit $663$ M\euro{} à $99{,}5\,\%$ ; moyenne de queue $1\,752{,}4$ M\euro{} \\
Intervalles & $\xi$ : bootstrap $[0{,}304\,;0{,}831]$, delta-méthode $[0{,}320\,;0{,}870]$ ;
  quantile $[411\,;1\,037]$, facteur $2{,}5$ \\
Hill & $1{,}32$ à $k=86$, $+120{,}8\,\%$ ; les six valeurs observées tombent dans l'intervalle à
  $90\,\%$ de $2\,000$ échantillons simulés sous $\xi=0{,}595$ \\
Seuil & $\hat\xi$ stable à $0{,}5\,\%$ près jusqu'au percentile $85$ ; de $0{,}60$ à $0{,}38$
  au-dessus (dans l'IC) ; $0{,}98$ en dessous (hors IC) ; règle de stabilité : $19{,}87$ contre
  $20{,}03$ ; règle permissive : $12{,}88$ et un quantile $+16{,}0\,\%$ \\
Tests & famille GPD : Anderson-Darling $p=0{,}96$, Kolmogorov-Smirnov $p=0{,}89$ ; paramètres
  publiés imposés : $0{,}99$ et $0{,}89$ ; couverture réelle de l'IC à $90\,\%$ : $86\,\%$ \\
Jackknife & un retrait à la fois : étendue $111{,}7$ M\euro{} ($17{,}9\,\%$ de l'IC) ; trois plus
  grosses pertes retirées : sous la borne basse ; cinq : $657{,}2\to327{,}3$, $\xi=0{,}22$ \\
\end{fichetab}
\end{parcoeur}

\begin{parcoeur}[À savoir par c\oe ur : la fréquence et le backtest]
\begin{fichetab}
Clé de répartition & $\lambda_{ref}=341=\sum_v\lambda_v$, $\lambda_v=\lambda_{ref}\,\pi_v$ ;
  elle n'alimente pas le moteur \\
Binomiale négative & $r=\lambda/(\varphi-1)$, $p=r/(r+\lambda)$ ; $\varphi$ gardé constant d'un
  scénario à l'autre, $r$ recalculé \\
Surdispersion & $1{,}19$ sur les $14\,944$ cellules firme-année ; $2{,}24$ $[1{,}24\,;5{,}18]$ sur
  la série détendancée (dix-huit points) ; $9{,}20$ posé au moteur ; $a=0{,}122$
  $[0{,}054\,;0{,}223]$ contre $0{,}60$ posé \\
Poisson & rejeté à $p\approx10^{-30}$ ; passer à la binomiale négative relève le SCR de
  $+9\,\%$ seulement ; quantile du nombre annuel $74$ (Gamma) contre $82$ (lognormal) \\
Backtest & origine glissante 2004--2025, $539$ incidents, douze années notées ; couverture
  $91{,}7\,\%$ (binomiale négative) contre $75{,}0\,\%$ (Poisson) ; forme de queue non rejetée,
  $p=0{,}4087$ \\
Dérive & échelle de queue $+4{,}00\,\%$ par an, dont la moitié d'inflation monétaire ; au moins
  $+24{,}7\,\%$ sur le quantile ; facteur entre états $3{,}344\to3{,}124$ ; écart $-1{,}9\,\%$ \\
Puissance & doubler le taux de dépassement se détecte à $80\,\%$ en $1\,811$ années à
  $99{,}5\,\%$, $905$ à $99\,\%$, $78$ à $90\,\%$ \\
\end{fichetab}
\end{parcoeur}

\begin{question}{\og Pourquoi les dépassements de seuil plutôt que les maxima par blocs ?\fg{}}
\rep{Les maxima par blocs ne gardent qu'une observation par bloc. Sur vingt-deux années, cela fait vingt-deux points. Les dépassements de seuil utilisent les $91$ pertes au-dessus de $20{,}03$~M\euro. Le théorème de Pickands, Balkema et de Haan assure que les deux méthodes mènent à la même famille de lois, avec le même indice de queue.}
\appui{\S~4.1, folio 31.}
\end{question}

\begin{question}{\og Écrivez le quantile de sévérité et donnez sa valeur.\fg{}}
\rep{$q_\alpha=u+\frac{\sigma}{\xi}\bigl[\bigl(\frac{1-\alpha}{\zeta_u}\bigr)^{-\xi}-1\bigr]$, valable pour $\xi$ entre $0$ et $1$. À $99{,}5\,\%$, il vaut $663$~M\euro{}. La moyenne de queue vaut $1\,752$. Il faut ajouter la phrase qui évite l'erreur : c'est le montant d'\emph{un} sinistre, pas un capital. Le capital est le quantile de la charge annuelle. Pour l'atteindre, il faut monter à environ $99{,}98\,\%$ par sinistre.}
\appui{\S~4.2, folios 32 et 33, encadré.}
\end{question}

\begin{question}{\og $\xi=0{,}595$ : votre moyenne est-elle infinie ?\fg{}}
\rep{Non. Entre $1/2$ et $1$, la variance est infinie, mais la moyenne existe. La formule du quantile et de la moyenne de queue suppose bien $\xi\in(0,1)$. La moyenne n'existerait plus qu'avec la calibration PRC, où $\xi$ vaut $1{,}03$.}
\appui{\S~4.1, folio 31 ; \S~4.2, folios 32 et 33.}
\end{question}

\begin{question}{\og Votre estimateur de Hill vaut $1{,}32$ : votre queue n'est-elle pas bien plus lourde ?\fg{}}
\rep{Si $\xi$ valait $1{,}32$, la moyenne n'existerait pas. Mais $1{,}32$ est exactement ce que Hill donne sur un petit échantillon quand $\xi=0{,}595$. J'ai simulé $2\,000$ échantillons avec le modèle publié. Les six valeurs observées tombent dans l'intervalle à $90\,\%$. De plus, Hill monte quand $k$ augmente, alors que $\hat\xi$ baisse quand le seuil monte. Si la queue valait vraiment $1{,}32$, ces deux tendances iraient dans l'autre sens.}
\appui{\S~4.5, folio 34.}
\end{question}

\begin{question}{\og L'indice de queue est-il stable quand on bouge le seuil ?\fg{}}
\rep{Non, et le mémoire ne dit pas le contraire. L'indice est stable près du seuil choisi. Quand on monte du percentile $85$ au $95$, il baisse de $0{,}60$ à $0{,}38$. Mais il reste dans l'intervalle déjà publié, donc il n'ajoute pas d'incertitude. Quand on descend, il remonte à $0{,}98$ et sort de l'intervalle. C'est le biais de seuil habituel, et c'est pourquoi je ne descends pas plus bas.}
\appui{\S~4.6, folios 35 et 36.}
\end{question}

\begin{question}{\og Pourquoi ce seuil-là ? Il vous arrange.\fg{}}
\rep{J'aime cette question, parce que la mesure va contre le soupçon. J'ai appliqué la règle de stabilité à dix seuils candidats. Elle retombe sur $19{,}87$~M\euro, contre $20{,}03$ publiés, avec le même quantile. C'est aussi le seuil qui s'ajuste le mieux. La règle la plus permissive descendrait à $12{,}88$ et donnerait un quantile \emph{plus haut} de $16\,\%$. Mon seuil ne maximise donc pas le chiffre du mémoire.}
\appui{\S~4.6, folios 35 et 36.}
\end{question}

\begin{question}{\og Avez-vous testé ce seuil avec Kolmogorov-Smirnov, et quels sont vos estimateurs avec leurs intervalles ?\fg{}}
\rep{Oui, après le dépôt, et je le présente comme un calcul complémentaire, hors rapport. Avec des paramètres estimés sur les mêmes données, la table usuelle du test est trop indulgente. J'ai donc calibré la $p$-value par bootstrap, en réestimant la loi à chaque tirage. Au seuil retenu, $D=0{,}043$ et $p=0{,}938$. Seuls les deux seuils les plus bas sont rejetés à $5\,\%$, par Anderson-Darling. Au seuil retenu, le maximum de vraisemblance donne $\xi=0{,}598$, intervalle à $95\,\%$ de $0{,}250$ à $0{,}890$. Les moments pondérés donnent $0{,}507$, Hill $1{,}377$. Je ne prétends pas que le test désigne le seuil optimal : à haut seuil, il reste peu d'excès, donc peu de puissance.}
\appui{Calcul complémentaire, script 119 (hors rapport) ; \S~4.6, folios 35 et 36.}
\end{question}

\begin{question}{\og $582$ pertes, $91$ excès : n'est-ce pas trop peu, et tout tient-il à trois sinistres ?\fg{}}
\rep{C'est peu, et c'est pour cela que la queue est la première source d'incertitude. Je la chiffre par un intervalle sur $\xi$ et je la propage jusqu'au capital. Retirer un seul sinistre ne fait jamais sortir le quantile de son intervalle. Retirer les trois plus gros le fait passer sous la borne basse. En retirer cinq le divise presque par deux. Ce n'est pas une fragilité, c'est une condition de lecture. À $99{,}5\,\%$, le quantile \emph{est} porté par les plus gros sinistres. Les retirer, c'est retirer ce qu'on mesure.}
\appui{\S~4.6, folios 35 et 36.}
\end{question}

\begin{question}{\og Votre surdispersion de $9{,}2$ est-elle mesurée ?\fg{}}
\rep{Non, elle est posée. L'indice de dispersion vaut $1+\lambda/r$, donc il change avec l'échelle. Il vaut $1{,}19$ sur les cellules firme-année. Il vaut $2{,}24$ sur la série annuelle sans tendance. Il vaut $9{,}20$ dans le moteur, par prudence. Ce sont trois niveaux de lecture : une échelle, une agrégation, puis une marge assumée.}
\appui{\S~4.3, folios 33 et 34.}
\end{question}

\begin{question}{\og Poisson ou binomiale négative, qu'est-ce que cela change au capital ?\fg{}}
\rep{Seulement $+9\,\%$. La moyenne ne change pas. La queue du nombre de sinistres change, mais à $\xi\approx0{,}6$ le quantile annuel est tiré par la plus grosse perte. La binomiale négative reste la bonne loi. Au backtest, elle couvre $91{,}7\,\%$ des années, contre $75{,}0\,\%$ pour Poisson. Le log-score, qui pénalise les intervalles trop larges, donne le même gagnant.}
\appui{\S~4.3, folios 33 et 34 ; \S~4.7, folio 36.}
\end{question}

\begin{question}{\og Votre intervalle à $90\,\%$ couvre-t-il vraiment $90\,\%$ ?\fg{}}
\rep{Non, il couvre $86\,\%$. Il est trop étroit, donc l'incertitude sur $\xi$ est sous-estimée. Il faudrait élargir la demi-largeur d'environ $10\,\%$. L'écart est petit, mais son sens compte. Il va du même côté que les deux autres défauts d'estimation : la sous-estimation.}
\appui{\S~4.6, folios 35 et 36.}
\end{question}

\begin{question}{\og On ne peut pas valider un quantile à $99{,}5\,\%$.\fg{}}
\rep{C'est exact, et le mémoire le chiffre au lieu de l'éviter. Pour détecter un taux de dépassement double du taux nominal, avec $80\,\%$ de puissance, il faudrait $1\,811$~années. Le backtest soutient la loi de comptage. Il ne rejette pas la forme de la queue. Mais il rejette la couverture des quantiles de sévérité : $5$ dépassements pour $1{,}4$ attendu à $99{,}5\,\%$, avec $p=0{,}0140$. Il ne valide donc pas le quantile réglementaire. Dire le contraire serait faux.}
\appui{\S~4.7, folio 36, encadré.}
\end{question}

\begin{question}{\og Votre backtest rejette le niveau. Pourquoi ne pas corriger ?\fg{}}
\rep{Il le rejette pour une raison que j'ai nommée et mesurée. L'échelle de sévérité dérive de $4{,}00\,\%$ par an dans la queue, dont la moitié vient de l'inflation. Cela donne au moins $+24{,}7\,\%$ sur le quantile. La calibration est gelée : je déclare et je chiffre l'écart au lieu de le corriger. Il ne gonfle pas ma thèse non plus. Avec la sévérité qui dérive, le facteur entre états baisse de $3{,}344$ à $3{,}124$. La bonne formule est « invariance d'échelle », pas « invariance à la dérive ».}
\appui{\S~4.7, folio 36, encadré.}
\end{question}

% ---------------------------------------------------------------------
\section*{Chapitre 5. Cascade dirigée entre piliers}
\addcontentsline{toc}{section}{Chapitre 5. Cascade dirigée entre piliers}
\reperes{folios 37 à 44. 5.1 Vasicek à un facteur (37), 5.2 généralisation~A (37), 5.3 généralisation~B (38), 5.4 Leontief (39), 5.5 branchement (40), 5.6 le seuil $K_j$ (41), 5.7 calibration (42), 5.8 formalisation (42), 5.9 vérifications (44).}

\enbref{La dépendance n'est pas ajoutée par une copule, elle est produite par la propagation d'un
incident le long d'une matrice dirigée ; la normalisation de Leontief garantit que la cascade
s'éteint, et trois propriétés qui portent l'argument sont vérifiées numériquement.}

\begin{parcoeur}[À savoir par c\oe ur : de Vasicek à la cascade]
\begin{fichetab}
Vasicek & $X_i=\sqrt\rho\,Y+\sqrt{1-\rho}\,\varepsilon_i$, défaut si $X_i\ge K_i$ ; deux
  défauts pour le cyber : dépendance symétrique (le jugement d'expert est dirigé à $81\,\%$) et
  probabilité de défaut non mesurable \\
Généralisation A & un facteur par pilier ; c'est le Vasicek multifacteur de Li (2000), pas une
  contribution ; $\rho_{ij}=\rho_j+u_{ij}$ (compromis hiérarchique) \\
$\rho_{ij}$ libre & jamais lisible sur un taux de conformité (la marge vaut $0{,}35$ quel que soit
  $\rho$) ; de l'ordre de $813$ périodes (approximation locale) pour séparer un écart de $0{,}10$ même à facteur connu ; $474$ au
  niveau du pilier ; $81$ sur une population entière ; $374$ entités sous le regroupement publié \\
Convention & $W_{jk}$ va de la source $j$ vers la cible $k$ : la ligne émet, la colonne reçoit ;
  P1 a la plus grande somme de ligne ($2{,}6$) et une somme de colonne trois fois plus petite que
  celle de P2 \\
Leontief & $W=g\,T/\max_j\sum_kT_{jk}$, donc
  $\rho(W)\le\max_j\sum_kW_{jk}\le g<1$ ; $\rho(T)=1{,}461$ et
  $(I-T)^{-1}$ a $21$ entrées négatives sur $25$ \\
Deux rayons & $\rho(W)=0{,}562\,g$, soit $0{,}506$ à $g=0{,}9$ ; $R_0=\rho(M)=0{,}054$, avec
  $M_{jk}=\Phi(\text{base}_j+W_{jk})-\Phi(\text{base}_j)$ et une incidence de fond de $4{,}5\,\%$ ;
  seuils critiques $g=1{,}78$ et $8{,}1$, hors domaine \\
Progéniture & $0{,}109$ (P1), $0{,}069$ (P4), $0{,}058$ (P2), $0{,}054$ (P3), $0{,}026$ (P5) :
  l'ordre de ROOT, par cohérence interne (même expert), pas par validation \\
\end{fichetab}
\end{parcoeur}

\begin{parcoeur}[À savoir par c\oe ur : les propriétés]
\begin{fichetab}
Descendance & $\mathbb E|S_a|\le[(I-W)^{-1}\mathbf 1]_a$ ; l'écart, de $10{,}6\,\%$ (P5) à
  $19{,}9\,\%$ (P1), compte les collisions \\
Loi exacte & $\Prob(S_a=A)=R(A)\,B(A)$, $B(A)=\prod_{j\in A,k\notin A}(1-W_{jk})$ \\
Moteur & une marche auto-évitante, même matrice moyenne que la cascade de la définition ; passer au
  branchement déplace le SCR de $-4\,\%$ \\
Monotonie & $W\le W'$ entraîne $S_a(W)$ dominé par $S_a(W')$ ; mais le capital n'est pas monotone
  en $A_{jk}$, d'où des sommets sans garantie \\
Ordre & la criticité est invariante par permutation si et seulement si $T$ est
  symétrique ; $22$ sous-ensembles sur $26$ changent de criticité \\
Co-occurrence & $2r\,S_{jk}+O(\lVert W\rVert^2)$ au premier ordre sous amorce uniforme ; la loi
  complète distingue $W$ de $W^\top$ (cascade seule) : $1{,}931$ contre $1{,}746$ pilier par sinistre,
  $62{,}31\,\%$ contre $50{,}77\,\%$ de multi-piliers \\
Transposer & capital $3{,}6\,\%$, perte moyenne $7{,}9\,\%$, priorité de P1 à P4 ; jusqu'à
  $1\,173$ M\euro{} sur l'ensemble admissible \\
Non-transitivité & réfutée : zéro violation sur $420$ triplets, valeur propre minimale
  $+0{,}405$, zéro cycle \\
Seuil $K_j$ & $K_j=F^{-1}(1-p_j)$, $p_j=\Prob(S\ge u_j)$, $u_j$ seuil de matérialité (art.~18) ;
  valeurs extrêmes : biais $-0{,}8\,\%$ contre $-63{,}9\,\%$ pour la lognormale \\
Calibration & probit, un logit gonflerait les magnitudes d'un facteur $2{,}1$ ; sur données
  simulées, corrélation $0{,}92$ et pente $1{,}08$ ; sur données réelles, rien \\
\end{fichetab}
\end{parcoeur}

\begin{question}{\og Votre cascade peut-elle exploser ?\fg{}}
\rep{Non, et c'est garanti par construction. La normalisation impose que chaque ligne de $W$ ait une somme d'au plus $g<1$. Donc $\rho(W)\le g<1$, et la série $I+W+W^2+\cdots$ converge. Le piège est documenté : la matrice brute de l'expert a un rayon spectral de $1{,}461$. De plus, $(I-T)^{-1}$ a $21$ entrées négatives sur $25$. Un choc positif donnerait alors un stress négatif ailleurs.}
\appui{\S~5.4, folios 39 et 40.}
\end{question}

\begin{question}{\og Démontrez la stabilité au tableau.\fg{}}
\rep{$W$ a des coefficients positifs. Chaque ligne a une somme d'au plus $g$, car $s_{\max}$ est la plus grande somme de ligne de $T$. Le rayon spectral est plus petit que n'importe quelle norme d'opérateur. Je prends la norme infinie, qui est le maximum des sommes de lignes. J'obtiens $\rho(W)\le\lVert W\rVert_\infty\le g<1$. La série de Neumann converge. Donc $(I-W)^{-1}$ existe, et ses coefficients sont positifs.}
\appui{annexe~A, folio 87.}
\end{question}

\begin{question}{\og Vous confondez deux rayons spectraux.\fg{}}
\rep{Non, je les distingue. $\rho(W)$ gouverne le système latent simultané. Il doit être inférieur à un pour que $(I-W)^{-1}$ existe. $R_0=\rho(M)$ gouverne la cascade d'incidents. Elle s'éteint si et seulement si $R_0<1$, comme une épidémie. Les valeurs sont $0{,}506$ et $0{,}054$. Les deux sont stables, mais pas au même niveau. La forme simultanée atteint son seuil plus de quatre fois plus vite ($1{,}78$ contre $8{,}1$).}
\appui{\S~5.5, folios 40 et 41, encadré.}
\end{question}

\begin{question}{\og Pourquoi un Vasicek, qui vient du risque de crédit ?\fg{}}
\rep{Parce que la structure est la même. Une variable latente franchit un seuil, et un facteur commun crée la dépendance. Je dis ce que j'emprunte : le Vasicek multifacteur à corrélation générale de Li (2000). Je l'utilise tel quel et je dis que ce n'est pas ma contribution. Ma contribution commence avec la contagion dirigée.}
\appui{\S~5.2, folios 37 et 38, encadré.}
\end{question}

\begin{question}{\og Pourquoi pas une sensibilité par entité et par pilier ?\fg{}}
\rep{Parce qu'une sensibilité ne se lit jamais sur un taux de conformité. Sa marge vaut $0{,}35$ quelle que soit la sensibilité. Elle se lit seulement dans la variation commune au fil du temps. Même si on connaissait le facteur, il faudrait environ $813$ périodes pour séparer un écart de $0{,}10$. Élargir le panel ne suffit pas : il faut encore $81$ périodes sur une population entière. C'est pourquoi je prends le compromis hiérarchique.}
\appui{\S~5.2, folios 37 et 38.}
\end{question}

\begin{question}{\og Pourquoi un seuil commun $K_j$ et pas un seuil $K_{ij}$ propre à chaque entité ?\fg{}}
\rep{Le cadre général s'écrit $K_{ij}=F_{ij}^{-1}(1-p_{ij})$, et il est correct. Mais $F_i$ seul suppose la même marge pour les cinq piliers, ce que la propagation ne garantit pas. Surtout, on ne peut pas l'estimer : aucune base n'observe la perte par pilier ni par entité, et il faudrait environ $813$ périodes pour séparer un écart de $0{,}10$. Je retiens donc une loi latente commune. Le capital ne dépend pas de ce choix : il vient de la cascade, d'une sévérité commune et de la descente d'échelle.}
\appui{Cours, section « Le seuil de déclenchement » ; \S~5.2, folios 37 et 38.}
\end{question}

\begin{question}{\og Votre moteur ne simule pas la cascade de votre définition.\fg{}}
\rep{C'est exact, et c'est écrit. Le moteur simule une marche qui ne repasse pas par un pilier déjà touché. La définition parle d'un processus de branchement. Les deux ont la même matrice moyenne et le même nombre moyen de descendants directs. Passer de l'un à l'autre change le capital de $-4\,\%$. Ce choix ne porte donc pas le résultat.}
\appui{\S~5.8, folios 42 à 44.}
\end{question}

\begin{question}{\og La progéniture retrouve exactement ROOT : c'est une validation ?\fg{}}
\rep{Non, c'est une cohérence interne. ROOT et TRANS viennent du même expert. Leur intérêt réel est de réduire les paramètres : ROOT se déduit de TRANS au lieu d'être posé. Cela fait cinq paramètres d'expert en moins.}
\appui{\S~5.5, folios 40 et 41, encadré.}
\end{question}

\begin{question}{\og Que représente la dépendance à l'ordre, au-delà d'une loi jointe des pertes ?\fg{}}
\rep{Le score d'une chaîne ordonnée dépend des liens de $T$ et du poids d'amorce. Sur les configurations examinées, $22$ des $26$ ensembles changent de criticité selon l'ordre. Un score portant seulement sur l'ensemble des piliers touchés ne représente pas cette différence. Une copule des seules pertes décrit une loi jointe ; il faudrait lui adjoindre une information sur les séquences pour comparer ces ordres. Le résultat ne prouve pas l'impossibilité générale de modèles à copules plus riches.}
\appui{\S~5.8 et 5.9, folios 42 à 44.}
\end{question}

\begin{question}{\og Pourquoi avoir abandonné la non-transitivité, qui donnait son titre au mémoire ?\fg{}}
\rep{Parce que je l'ai testée de trois façons, et aucune ne tient. Il y a zéro violation de transitivité sur $420$ triplets. La matrice de corrélation est parfaitement valide (plus petite valeur propre $+0{,}405$). Il y a zéro cycle de dominance. Il valait mieux la réfuter que la garder. Un travail qui réfute sa propre hypothèse de départ est plus crédible.}
\appui{\S~5.9, folio 44.}
\end{question}

\begin{question}{\og Que coûte de se tromper de sens ?\fg{}}
\rep{Si on transpose la matrice, le capital change de $3{,}6\,\%$ et la perte moyenne de $7{,}9\,\%$. Surtout, la priorité de remédiation passe de P1 à P4. Sur l'ensemble admissible, l'écart entre une matrice et sa transposée atteint $1\,173$~M\euro. Les statistiques disponibles ne les distinguent pas, mais l'effet est réel.}
\appui{\S~5.9, folio 44.}
\end{question}

\begin{question}{\og Pourquoi un probit plutôt qu'un logit ?\fg{}}
\rep{Parce que le modèle générateur a un bruit gaussien et un seuil. L'estimateur est donc un probit. Un logit gonflerait les magnitudes d'un facteur $2{,}1$. Sur des données simulées, le probit retrouve $W$ (corrélation $0{,}92$, pente $1{,}08$). Ce qui échoue sur les vraies données, c'est la donnée, pas la méthode.}
\appui{\S~5.7, folio 42.}
\end{question}

\lepiege{L'argument \og la co-occurrence ne voit que $S$\fg{} ne vaut qu'au premier ordre et sous
amorce uniforme. La loi complète des piliers touchés distingue $W$ de $W^\top$ (fourche contre
collisionneur). Ne jamais présenter la frontière comme une impossibilité de principe.}

% ---------------------------------------------------------------------
\section*{Chapitre 6. Modélisation multi-états de la conformité}
\addcontentsline{toc}{section}{Chapitre 6. Modélisation multi-états de la conformité}
\reperes{folios 45 à 49. 6.1 deux échelles de temps (45), 6.2 Vasicek multi-états (45), 6.3 les quatre canaux (46), 6.4 les valeurs de $g$ (47), 6.5 l'état par pilier (48), 6.6 Markov (49), 6.7 le couplage (49).}

\enbref{L'état de conformité n'est ni binaire ni figé : un Vasicek à trois états le fixe à
l'échelle rapide, une chaîne de Markov le fait évoluer à l'échelle lente, et il pilote le capital
par quatre canaux, sans aucune pénalité ajoutée.}

\begin{parcoeur}
\begin{fichetab}
Latente & $C^\star=\gamma\,\Theta+\sqrt{1-\gamma^2}\,\varepsilon$, $\gamma=0{,}68$, donc une
  corrélation $\gamma^2=0{,}462$ entre deux piliers \\
Seuils & $K_{\text{bas}}=\Phi^{-1}(p_{\text{NC}})$,
  $K_{\text{haut}}=\Phi^{-1}(p_{\text{NC}}+p_{\text{PC}})$ ; état actuel posé NC $35\,\%$, PC
  $35\,\%$, C $30\,\%$ \\
Quatre canaux & fréquence $21{,}56\to53{,}62$ incidents par an ; détection, $p_u$ multiplié par
  $0{,}85$ ou $1{,}20$ ($0{,}128\to0{,}181$) ; propagation $g$ de $0{,}45$ à $0{,}90$ ;
  accumulation P4 de $0$ à $0{,}68$ \\
Multiplicateur & $\lambda_{\text{NC}}=\lambda_{\text{C}}\sum_v\pi_v m_v=21{,}56\times2{,}4875$ ;
  centres $2{,}40$, $2{,}30$, $2{,}25$, $5{,}50$, $1{,}45$ ; aux bornes, de $1{,}933$ à $3{,}042$,
  soit $41{,}7$ à $65{,}6$ incidents ; les identifiants font $6{,}3\,\%$ des incidents et
  $13{,}9\,\%$ du multiplicateur \\
Valeurs de $g$ & $0{,}45$, $0{,}68$, $0{,}90$, posées ; lemme de croissance sur vingt points à
  nombres aléatoires communs ; entité $116{,}4$, $143{,}2$, $169{,}0$ M\euro{} ($+45{,}2\,\%$) ;
  correspondance très étroite $+3{,}9\,\%$ ; forfait $450$ M\euro{} plat \\
Élasticité & $0{,}41$ sur $g\in[0{,}20\,;0{,}95]$ ; un écart de plus de $10\,\%$ dès que $g$
  passe de $0{,}90$ à $0{,}75$ \\
Dépendance entre états & surcroît $\delta$ admissible jusqu'à $0{,}38$ ; $\Prob(\text{tous NC})$
  de $0{,}068$ à $0{,}088$ ; capital espéré $+2$ M\euro{} ($0{,}02\,\%$) ; forme additive
  $R^2=0{,}9945$, contributions $3\,524$ (P1), $2\,863$ (P4), $2\,022$ (P2), $1\,577$ (P3), $897$
  (P5) \\
Markov & NC vers PC vers C, C absorbant ; séjours Erlang-$k$ (mode strictement positif) ;
  $\overline{\mathrm{SCR}}(t)=\sum_e\Prob(e,t)\,\mathrm{SCR}_e$, une moyenne de quantiles (écart
  de $9{,}3\,\%$ au quantile du mélange) \\
Couplage & un seul $\Theta$ : $\mu(\Theta)=\mu_0e^{\beta_{\text{rem}}\Theta}$ et
  $\mathrm{SCR}_e(\Theta)=\mathrm{SCR}_e\,e^{-\beta_{\text{casc}}\Theta}$ \\
\end{fichetab}
\end{parcoeur}

\begin{question}{\og Qu'est-ce qui vous dit que la conformité réduit la propagation ?\fg{}}
\rep{Rien dans les données, et c'est écrit. C'est la seule hypothèse qualitative qui reste. Le règlement prend lui-même parti, puisqu'il impose ces exigences. Mais ma thèse n'en dépend pas. Le capital augmente avec $g$. Toute correspondance qui respecte l'ordre redonne donc l'écart. L'élasticité d'arc vaut $0{,}41$. Un écart de plus de $10\,\%$ demande seulement un écartement de $0{,}15$ sur $g$. On demande un ordre de grandeur, pas une mesure.}
\appui{\S~6.4, folios 47 et 48, encadré.}
\end{question}

\begin{question}{\og Vos trois valeurs de $g$ sont arbitraires.\fg{}}
\rep{Oui, elles sont posées, parce qu'aucune entité ne publie son état de conformité. Voici ce qui tient sans elles. Le capital augmente avec $g$ : je l'ai vérifié sur vingt points, sans aucune baisse. L'existence et le sens de l'écart ne dépendent donc pas de $g$. Seule l'amplitude est un scénario. Même avec une correspondance très étroite, de $0{,}85$ à $0{,}90$, l'écart reste de $+3{,}9\,\%$ quand le forfait reste plat.}
\appui{\S~6.4, folios 47 et 48.}
\end{question}

\begin{question}{\og D'où viennent vos multiplicateurs de fréquence ?\fg{}}
\rep{Ils se construisent vecteur par vecteur. Pour chaque vecteur, je prends sa part dans Hackmageddon. Je la multiplie par le centre d'une fourchette tirée d'un effet publié : la formation contre l'hameçonnage, l'authentification forte, etc. Le total vaut $2{,}4875$, soit $53{,}62$ incidents par an. Le multiplicateur est semi-ancré : sa provenance est tracée, mais sa valeur est posée. Aux bornes des fourchettes, la fréquence non conforme va de $41{,}7$ à $65{,}6$.}
\appui{\S~6.3, folios 46 et 47.}
\end{question}

\begin{question}{\og Être conforme aux piliers 1 et 2 aide pour le 3 : vous l'ignorez ?\fg{}}
\rep{Je ne l'ignore pas. La dépendance existe déjà : un facteur commun de charge $0{,}68$, soit une corrélation de $0{,}462$ entre deux piliers. Ce qui manque, c'est sa structure, qui est la même pour toutes les paires. J'ai ajouté un surcroît de corrélation sur les seules paires concernées. La loi des configurations bouge nettement. Mais le capital espéré ne bouge que de $2$~M\euro. C'est structurel : le capital des $32$ configurations est presque additif ($R^2=0{,}9945$). Et l'espérance d'une fonction additive ne dépend que des marges. Réserve : cet argument vaut pour l'espérance, pas pour un quantile.}
\appui{\S~6.5, folios 48 et 49.}
\end{question}

\begin{question}{\og D'où sort la corrélation de $0{,}462$ ?\fg{}}
\rep{Elle n'est pas posée à part. C'est le carré de la charge du facteur commun : $\gamma^2=0{,}68^2$.}
\appui{\S~6.2 et 6.5, folios 45 à 49.}
\end{question}

\begin{question}{\og Pourquoi des séjours de type Erlang plutôt qu'exponentiels ?\fg{}}
\rep{Une exponentielle a sa densité maximale en zéro. Elle permettrait une remédiation presque instantanée. Un projet DORA a une durée minimale. Une Erlang la représente, car son mode est strictement positif. Ce n'est pas la même Erlang que la loi qui mélange la fréquence. Pour celle-ci, la forme estimée est $0{,}63$, qui n'est pas entière.}
\appui{\S~6.6, folio 49.}
\end{question}

\begin{question}{\og Votre trajectoire de capital est-elle un quantile ?\fg{}}
\rep{Non. C'est une moyenne de quantiles, pondérée par la loi de l'état à chaque horizon. Le quantile du mélange serait autre chose. La queue d'un mélange vient de ses composantes les plus lourdes. L'écart est de $9{,}3\,\%$ sur les $32$ configurations. J'ai déclaré cette convention et je la tiens dans tout le mémoire.}
\appui{\S~6.6, folio 49.}
\end{question}

\begin{question}{\og Pourquoi deux échelles de temps ?\fg{}}
\rep{Parce que la propagation dans un sinistre et l'évolution de la conformité d'une année à l'autre n'ont pas la même horloge. Les mélanger reviendrait à expliquer un enchaînement de quelques heures par une tendance longue. C'est l'architecture classique de la migration de rating. Le préprint climatique sépare aussi ses indices, pour la même raison.}
\appui{\S~6.1, folio 45 ; \S~6.7, folio 49.}
\end{question}

% ---------------------------------------------------------------------
\section*{Chapitre 7. Identifiabilité de la direction de contagion}
\addcontentsline{toc}{section}{Chapitre 7. Identifiabilité de la direction de contagion}
\reperes{folios 50 à 55. 7.2 trois échecs (50), 7.3 réversibilité (51), 7.4 MOVEit (52), 7.5 corpus (52), 7.6 corpus étendu (53), 7.7 biais de narration (54), 7.8 séquences (54), 7.9 bayésien (55), 7.10 frontière (50).}

\enbref{La direction de la contagion n'est pas identifiable sur les données publiques agrégées,
quand sa co-occurrence l'est : trois sources échouent pour trois raisons distinctes, la
réversibilité en donne la raison de principe, et un corpus de rapports d'enquête corrobore la
structure sans calibrer les coefficients.}

\begin{parcoeur}[À savoir par c\oe ur : les trois échecs]
\begin{fichetab}
Il faudrait & des co-occurrences dans une même entité, une antériorité fiable, une taxonomie par
  pilier \\
PRC & $74\,783$ incidents, $25\,009$ organisations ; catégorie indéterminée $52{,}7\,\%$ ;
  survenance exploitable dans $49{,}6\,\%$ ; $71$ transitions ($3{,}6$ par cellule), $34$ en
  finance : pas de taxonomie, pas de matière \\
OpRisk & $4\,329$ transitions annuelles (environ $103$ par cellule) ;
  $T(M)=\tfrac12\lVert M-M^\top\rVert_1=119$ contre $128\pm27$ au placebo, $z=-0{,}33$ ; aucune
  des $18$ paires significative : le pas annuel efface l'ordre \\
VERIS & $10\,591$ incidents ($980$ en finance) ; \texttt{control\_failure} renseigné $13$ fois
  ($0{,}12\,\%$), chronologie datée sur deux étapes $126$ fois ($1{,}19\,\%$), sur trois $10$
  fois ; les deux ensemble : $1$ incident, aucun en finance \\
Temps inversé & $M_{\text{inv}}=M^\top$ par construction : corrélation $-1$ quel que soit le
  contenu, le test est dégénéré \\
Réversibilité & courant $J_{jk}=\pi_jP_{jk}-\pi_kP_{kj}$ ; la forme de Dirichlet ne voit que la
  partie réversible ; production d'entropie $\sigma\ge0$, nulle si et seulement si réversible \\
MOVEit & $188$ organisations financières, fenêtres de $180$ jours ; différence de différences
  $+1{,}17$ $[+1{,}07\,;+1{,}27]$ ; placebo un an plus tôt $+1{,}11$ ; différence descriptive $+0{,}06$, sans test du contraste \\
\end{fichetab}
\end{parcoeur}

\begin{parcoeur}[À savoir par c\oe ur : le corpus et la frontière]
\begin{fichetab}
Sept rapports & CSRB, FCA et PRA, GAO, SEC, OCC ; $17$ transitions ; $z=+3{,}93$
  ($p=0{,}0004$) ; P1 : $9$ sorties, $0$ entrée ; P2 : $0$ sortie, $7$ entrées \\
Dix rapports & plus CrowdStrike 2024, Log4Shell, Raphaels Bank 2015 ; $22$ transitions,
  $z=+5{,}12$ ; sources primaires seules $+4{,}08$ ; perspective entité seule $+3{,}62$ \\
Matrice ordonnée & $p_{jk}\sim\mathrm{Beta}(1+m_{jk},1+N_j-m_{jk})$ ; $12$ liens sur $20$
  crédiblement sous $0{,}5$, aucun au-dessus ; P3 vers P2 $0{,}77$ $[0{,}48\,;0{,}95]$, P1 vers P4
  $0{,}68$ $[0{,}40\,;0{,}89]$ ; les quatre sorties de P2 et trois des quatre entrées de P1 \\
Avec l'expert & corrélation non significative, $+0{,}08$ (matrice) et $+0{,}31$
  (antisymétrique, $p=0{,}19$) ; dénominateur corrigé : niveau $0{,}27\to0{,}07$, part de
  direction $60\,\%$ inchangée \\
Biais de narration & $11$ des $22$ arêtes partent de P1 ; $p$ exacte $1{,}5\cdot10^{-5}$ ;
  quatre arêtes retournées au pire : $0{,}084$ ; les onze retirées : $0{,}0039$ ; les sept
  rapports retirés : $0{,}50$ \\
En capital & P1 muette : entité $169{,}0\to127{,}1$ M\euro{} ($-24{,}8\,\%$), soit $0{,}81$ fois
  la bande de $51{,}7$ ; priorité de P1 à P4 dès $a=0{,}50$ ; écart entre états
  $+45{,}2\,\%\to+29{,}1\,\%$ \\
Bayésien & $A=S(2\theta-1)$, $\theta\sim\mathrm{Beta}(\alpha_0+n_{lo},\beta_0+n_{hi})$ ;
  crédible $81\,\%$ plus resserré que les bornes ; deux paires sur six tranchées \\
Trois niveaux & structure revendiquée ; direction paire par paire non revendiquée ; amplitude non
  identifiable (une bande de $656$ M\euro{} même les dix sens connus) \\
\end{fichetab}
\end{parcoeur}

\begin{question}{\og Pourquoi la direction n'est-elle pas identifiable ?\fg{}}
\rep{Parce que trois sources échouent, chacune pour sa raison. La chronologie PRC n'a ni la taxonomie ni la matière. OpRisk a le volume ($103$ observations par cellule), mais son pas annuel efface l'ordre. VERIS a les bons champs, mais ne les remplit presque jamais : une fois sur $10\,591$. Ce n'est donc pas un manque de volume. La théorie des processus réversibles explique le reste : une statistique qui ne change pas quand on inverse le temps ne peut pas voir un courant.}
\appui{\S~7.2, folios 50 et 51 ; \S~7.3, folios 51 et 52.}
\end{question}

\begin{question}{\og C'est donc une impossibilité de principe ?\fg{}}
\rep{Non, et le mémoire le dit clairement. Au premier ordre et avec une amorce uniforme, la co-occurrence ne dépend que de $S$. Mais la loi complète des piliers touchés par incident distingue $W$ de $W^\top$, par exemple une fourche d'un collisionneur. La frontière est celle des statistiques que les sources donnent. D'où la quatrième exigence du registre : noter tous les piliers qui ont cédé.}
\appui{\S~7.3, folios 51 et 52, encadré ; \S~8.1 (proposition de non-identifiabilité), folio 56.}
\end{question}

\begin{question}{\og Le test du temps inversé ne suffisait-il pas ?\fg{}}
\rep{Il est dégénéré sur des comptages. Inverser la chronologie donne la transposée par construction, donc une corrélation de $-1$ quel que soit le contenu. Il ne prouve rien. Seul le placebo est valide : on permute les années à l'intérieur de chaque firme.}
\appui{\S~7.2, folios 50 et 51, encadré.}
\end{question}

\begin{question}{\og Votre expérience naturelle MOVEit, qu'a-t-elle donné ?\fg{}}
\rep{Rien, et c'est un résultat. La différence de différences brute était grande et significative : $+1{,}17$ jour-brèche, avec un intervalle qui exclut zéro. Mais le même dispositif, calé sur une fausse date un an plus tôt, donnait $+1{,}11$. L'effet venait donc du choix du groupe de contrôle. Une fois le placebo retiré, il reste $+0{,}06$. De plus, la cible serait restée non attribuable, car la taxonomie est un vecteur d'attaque.}
\appui{\S~7.4, folio 52.}
\end{question}

\begin{question}{\og Vous parlez de martingale : êtes-vous en mesure risque-neutre ?\fg{}}
\rep{Non. J'emploie le mot « martingale » au sens des chaînes de Markov. La forme de Dirichlet est la variation quadratique de la martingale associée à une fonction test, sous la mesure physique. Il n'y a ni mesure risque-neutre, ni actif répliqué, ni test de martingalité comme pour les générateurs de scénarios. Le mémoire ne valorise rien.}
\appui{\S~7.3, folios 51 et 52.}
\end{question}

\begin{question}{\og Que prouvent vos rapports d'enquête ?\fg{}}
\rep{Ils prouvent la structure : P1 en amont et P2 en aval, avec $z=+3{,}93$ puis $+5{,}12$. Le signal tient avec les sources primaires seules, et avec la seule perspective entité. Ils ne prouvent pas les coefficients : la corrélation avec la matrice d'expert n'est pas significative. Le corpus montre bien où la contagion ne va pas. Il ne montre pas où elle est forte.}
\appui{\S~7.5 et 7.6, folios 52 à 54.}
\end{question}

\begin{question}{\og Votre corpus sur-attribue à la gouvernance, et vous concluez qu'elle domine.\fg{}}
\rep{C'est l'objection la plus juste, et le mémoire la contient avant que vous ne la posiez. Une enquête officielle a pour mission de trouver une cause racine. Mais je mesure ce biais. Dans le pire cas, il faudrait avoir retourné quatre des onze arêtes de P1 pour perdre la significativité. Même en les retirant toutes, la direction reste significative, portée par P3 et P4 vers P2. En capital, rendre P1 muette ne coûte que $0{,}81$ fois la bande déjà publiée. Ce qui bouge, c'est la priorité : elle passe à P4 dès qu'on réduit la gouvernance de moitié.}
\appui{\S~7.7, folio 54.}
\end{question}

\begin{question}{\og Un seul codeur a lu vos rapports.\fg{}}
\rep{Oui, c'est une limite déclarée. Le second codage en aveugle est préparé. Le script d'accord ne produit rien sans le second jeu. En attendant, je borne le biais de narration par le pire cas : donner raison au critique sur tout coûte moins que l'ignorance de la direction déjà publiée. Un second codage ajouterait une mesure d'accord entre codeurs. Il ne corrigerait pas le niveau.}
\appui{\S~7.7, folio 54 ; supplément S2, \S~S2.3.}
\end{question}

\begin{question}{\og Pourquoi pas un bayésien, qui donnerait une loi plutôt qu'une bande ?\fg{}}
\rep{Le mémoire le fait. L'intervalle crédible à $90\,\%$ est $81\,\%$ plus serré que les bornes. Mais seules deux paires documentées sur six sont tranchées. Et l'intervalle crédible dépend d'un prior, alors que les bornes n'en dépendent pas. Je publie les deux. Les deux objets sont conditionnels aux hypothèses de construction ; l'intervalle crédible ajoute un prior sur la direction.}
\appui{\S~7.9, folio 55.}
\end{question}

\begin{question}{\og Et les triplets ordonnés, qui portent plus d'information que les paires ?\fg{}}
\rep{Je les ai explorés, puis fermés, pour trois raisons. La comparaison naturelle n'a aucune donnée : P2 n'émet jamais, et P1 n'est jamais atteint. Le test de dépendance au pilier précédent n'a aucune puissance. Et l'absence de cycle s'obtiendrait deux fois sur trois par hasard. Il faudrait de $158$ à $589$ chemins passant par un même pilier, contre neuf aujourd'hui.}
\appui{\S~7.8, folios 54 et 55.}
\end{question}

\begin{question}{\og Vos praticiens vous contredisent.\fg{}}
\rep{Deux praticiens sur trois contestent le sens unique de deux arcs : la gouvernance comme source pure, et les incidents comme puits. La plupart de leurs exemples décrivent une rétroaction corrective. Elle relève de l'évolution des états de conformité, pas de la propagation d'une perte. Et comme l'ensemble admissible ne fixe aucune direction, le capital publié ne bouge pas. Surtout, aucun ne distingue un ordre d'une cause commune. Ils confirment donc la frontière, pas la matrice.}
\appui{\S~7.5, folios 52 et 53, encadré ; supplément S4.}
\end{question}

\begin{question}{\og Qu'est-ce que votre frontière en trois niveaux ?\fg{}}
\rep{La structure est revendiquée : P1 en amont, P2 en aval, avec son facteur de confusion déclaré. La direction paire par paire ne l'est pas. Seules deux paires sur six sont tranchées, donc les bornes restent. L'amplitude n'est pas identifiable non plus par cette voie. Même avec les dix sens connus, il resterait une bande de $656$~M\euro. Je revendique un ordre partiel, pas une matrice.}
\appui{\S~7.10, folios 50 à 55.}
\end{question}

% ---------------------------------------------------------------------
\section*{Chapitre 8. Identification partielle et bornes de capital}
\addcontentsline{toc}{section}{Chapitre 8. Identification partielle et bornes de capital}
\reperes{folios 56 à 59. 8.1 le modèle en couches (56), 8.2 ce que le capital devient (57), 8.3 poser ou borner (58), 8.4 le niveau est illustratif (58), 8.5 l'échelle de repli (58).}

\enbref{Puisque la direction n'est pas identifiée, on ne la pose pas : on énumère les $1\,024$
matrices extrêmes compatibles avec la co-occurrence et l'on publie le capital en bornes, le socle
sans contagion restant défendable sans aucun paramètre non identifié.}

\begin{parcoeur}
\begin{fichetab}
Socle & $W=0$ : $5\,275$ M\euro{}, le plancher défendable \\
Ensemble admissible & $\mathcal W(t)=\{S+A : A^\top=-A,\ |A_{jk}|\le t\,S_{jk},\
  \rho(S+A)<1\}$ ; la contrainte $|A|\le S$ vient de la seule positivité ; $t=1$ est l'ignorance
  totale \\
Bornes & $2^{10}=1\,024$ sommets ; $[6\,858\,;8\,697]$ M\euro{}, facteur $1{,}27$, soit $+30$ à
  $+65\,\%$ au-dessus du socle ; le point d'expert, $8\,110$, est un point de l'ensemble \\
Sommets & l'atteinte des extrêmes n'est pas garantie (monter $A_{jk}$ baisse $W_{kj}$) ; aucun
  des $3\,000$ points intérieurs ne sort des bornes, en VaR ni en perte moyenne \\
Perte moyenne & bornée dans $[849\,;1\,100]$ M\euro{}, moins bruitée que la VaR \\
Valeur de $t$ & passer de $t=1$ à $t=0{,}5$ fait tomber la largeur de $1\,839$ à $695$
  ($-62\,\%$) \\
Poser ou borner & poser $W$ : un point, $8\,110$, qu'un panel biaisé déplace de $777$ M\euro{}
  ($10\,\%$) sans signal ; borner : la bande ne bouge pas, $79\,\%$ du capital est insensible au sens (socle $61\,\%$,
  plancher commun aux sommets $18\,\%$) \\
Dominance & P1 en tête dans $45\,\%$ des configurations, P4 dans $36\,\%$ ($81\,\%$ à eux deux),
  P3 dans moins de $10\,\%$, P5 jamais \\
Niveau et rapport & le niveau varie d'un facteur $41$ (socle PRC $239$, fréquence triplée
  $11\,709$) ; les rapports d'au plus $33\,\%$ ; le groupe de tête est le même \\
Repli et cran nul & la famille à un facteur réduit l'étendue de $25\,\%$ sur 400 tirages ; à $u=0$,
  gains de P1 et P4 proches ($1\,051$ et $1\,052$ M\euro), $353$ d'écart au classeur, mais $S$ change
  et une asymétrie subsiste \\
\end{fichetab}
\end{parcoeur}

\begin{question}{\og Alors, quel est votre SCR ?\fg{}}
\rep{Le niveau n'est pas une mesure, et le mémoire le démontre au lieu de s'en excuser. Il varie d'un facteur $41$ quand on change seulement la source de sévérité et la fréquence. Ce que le modèle détermine, c'est un intervalle et son prix : entre $6\,858$ et $8\,697$~M\euro{} sur les $1\,024$ sommets. La largeur de $1\,839$ est calculée aux sommets, sous les hypothèses retenues. Ce qui résiste au changement de source, c'est le rapport au socle et le groupe de tête des piliers.}
\appui{\S~8.2, folio 57 ; \S~8.4, folio 58.}
\end{question}

\begin{question}{\og Pourquoi borner, au lieu d'estimer ?\fg{}}
\rep{La direction n'est pas identifiée par les sources examinées. Ajouter des données du même type ne suffit donc pas nécessairement. Nous fixons la partie symétrique et faisons varier les asymétries dans un ensemble déclaré. Le paramètre $t$ règle linéairement l'amplitude admissible de chaque coefficient ; la largeur de la plage de capital doit être calculée, elle n'est pas démontrée linéaire en $t$.}
\appui{\S~8.1 et 8.2, folios 56 et 57 ; chapitre~7, folios 50 à 55.}
\end{question}

\begin{question}{\og Qu'est-ce qui garantit que les sommets donnent les extrêmes ?\fg{}}
\rep{Rien ne le garantit a priori, et le mémoire le dit. Monter $A_{jk}$ augmente $W_{jk}$ mais baisse $W_{kj}$. La réponse n'est donc pas monotone. Je procède par réfutation : je tire $3\,000$ points intérieurs de façon uniforme. Aucun ne sort des bornes, ni en VaR ni en perte moyenne. C'est une hypothèse vérifiée par le calcul, pas un théorème, et c'est écrit ainsi.}
\appui{\S~8.2, folio 57, figure 8.1.}
\end{question}

\begin{question}{\og Que vaut le point d'expert, $8\,110$ ?\fg{}}
\rep{C'est un point de l'ensemble, jamais sa mesure. Si le panel qui le fixe est biaisé, il se déplace de $777$~M\euro. La priorité change alors de pilier, et rien ne prévient le lecteur. La bande, elle, contient tous ces écarts et ne bouge pas d'un euro.}
\appui{\S~8.3, folio 58.}
\end{question}

\begin{question}{\og À quoi servirait alors une élicitation ?\fg{}}
\rep{Elle pourrait préciser les directions et départager P1 et P4, qui concentrent $81\,\%$ des matrices tirées dans ce diagnostic. P5 n'arrive jamais premier dans cet échantillon, mais ce n'est pas une exclusion démontrée sur tout l'ensemble admissible. Les fréquences de priorité doivent rester rattachées aux matrices effectivement testées.}
\appui{\S~8.3, folio 58.}
\end{question}

\begin{question}{\og Pourquoi regarder la perte moyenne et pas seulement la VaR ?\fg{}}
\rep{Avec une queue lourde, la VaR dépend surtout de la plus grosse perte. Elle réagit peu à la structure de la cascade. Un tableau de sensibilité fondé sur elle serait plat, et donc faussement rassurant. La perte moyenne, bornée entre $849$ et $1\,100$, est bien moins bruitée. C'est elle qui permet de lire la décision, c'est-à-dire l'ordre de remédiation.}
\appui{\S~8.2, folio 57, encadré.}
\end{question}

\begin{question}{\og Si la direction pèse si peu sur le niveau, à quoi sert-elle ?\fg{}}
\rep{Elle éclaire la décision : à $S$ fixé, la direction peut modifier la priorité entre piliers. Au cran nul, les gains de P1 et P4 sont proches, $1\,051$ contre $1\,052$~M\euro{}, pour $353$~M\euro{} d'écart au classeur. Mais $S$ change et une asymétrie subsiste : ce contraste n'isole pas la direction.}
\appui{\S~8.5, folios 58 et 59, encadré.}
\end{question}

\begin{question}{\og Le niveau est-il au moins une indication ?\fg{}}
\rep{C'est une illustration, pas une mesure. Le niveau varie d'un facteur $41$ quand on change des \emph{entrées} dont le périmètre n'est pas transposable. Les rapports au socle varient d'au plus $33\,\%$, soit trente fois moins. Le groupe de tête reste le même. Le rapport le moins stable, la largeur de bande divisée par le socle, varie de $32{,}9\,\%$. Le relatif n'est donc pas parfaitement stable, mais il varie bien moins que le niveau.}
\appui{\S~8.4, folio 58.}
\end{question}

% ---------------------------------------------------------------------
\section*{Chapitre 9. Résultats}
\addcontentsline{toc}{section}{Chapitre 9. Résultats}
\reperes{folios 60 à 70. 9.1 SCR par état et chiffre de tête (60), 9.2 décomposition par pilier (62), 9.3 Shapley et Euler (63), 9.4 trajectoire (65), 9.5 priorisation (65), 9.6 robustesse (65), 9.7 sensibilité à la propagation (65), 9.8 copule (66), 9.9 regard réglementaire (66), 9.10 ordre de grandeur (67), 9.11 entités réelles (68), 9.12 agrégation (69), 9.13 détenir ou transférer (70).}

\enbref{Le combien : sur les quatre canaux, le besoin de capital du secteur passe de $6\,049$ à
$20\,188$~M\euro{} ; le chapitre l'attribue, le déroule dans le temps, le confronte à une copule, à
la Formule Standard et à quatre entités réelles, puis le traduit en décision.}

\begin{parcoeur}[À savoir par c\oe ur : le chiffre de tête]
\begin{fichetab}
Quatre lectures & score seul, $g$ figé au niveau NC : $7\,987\to16\,847$ ($\times2{,}11$) ;
  fréquence et propagation : $6\,664\to15\,074$ ($\times2{,}26$) ; plus détection :
  $5\,932\to16\,595$ ($\times2{,}80$) ; quatre canaux : $6\,049\to20\,188$ ($\times3{,}34$) \\
Pourquoi la quatrième & la seule où l'état conforme l'est sur tous les canaux du modèle ; les
  trois autres sont des remédiations partielles, publiées comme telles \\
Détection & résolue : Shapley $2\,928\pm343$, fermeture $4\,562\pm819$, signe stable sur seize
  graines \\
Interaction & canaux isolés $9\,138$, interaction $5\,001$ ($35\,\%$) ; produit des quatre facteurs
  isolés $3{,}47$ contre $3{,}34$ publié : une composition presque multiplicative \\
Fréquence & seule, $10\,377$ ($\times1{,}72$) ; les trois autres réunis, $11\,413$ ($\times1{,}89$) ;
  aucun autre canal seul n'atteint $8\,000$ \\
Réconciliation & conforme autour de $5\,900$ (lecture~C, $5\,932$) contre socle $5\,275$ à $W=0$ :
  deux points d'un même modèle \\
Sources & le même état conforme vaut $8\,301$ sous OpRisk et $2\,589$ sous PRC (lecture~A) \\
Surcoût (lecture B) & NC contre C, OpRisk : médiane $7\,401$, IC à $90\,\%$ $[2\,830\,;38\,793]$,
  $100\,\%$ des tirages positifs \\
Mesure & $\mathrm{VaR}_{99{,}5\,\%}(L)$ \emph{non centrée} : la perte attendue n'étant couverte
  par aucune provision, elle n'est pas déduite \\
\end{fichetab}
\end{parcoeur}

\begin{parcoeur}[À savoir par c\oe ur : attribution, décision, comparaison]
\begin{fichetab}
Entre piliers & interaction de signe non résolu en VaR ($-137$ à $+1\,395$ sur six graines),
  positive en perte moyenne ($+532$ à $+659$) ; à ne jamais ajouter à celle entre canaux \\
Shapley & $\varphi_j=\sum_S\frac{|S|!(5-|S|-1)!}{5!}[v(S\cup\{j\})-v(S)]$ ; P1 $34{,}0\,\%$, P4
  $25{,}7\,\%$, P2 $17{,}1\,\%$, P3 $13{,}6\,\%$, P5 $9{,}6\,\%$ ; somme $10\,927$, marginaux
  $9\,531$ ; amorces ROOT $30$, $27$, $18$, $15$, $9\,\%$ \\
Euler & $C_j=\mathbb E[L_j\mid L\ge\mathrm{VaR}]$ ; presque plat sous PRC (P2 et P4 $23\,\%$, P1
  $20$, P3 $18$, P5 $17$) ; indicatif sous OpRisk (variance infinie) \\
Queue & $\mathrm{TVaR}/\mathrm{VaR}\approx2{,}1$ contre l'asymptote $1/(1-\xi)=2{,}5$ \\
Trajectoire & lecture~B : $10\,752\to7\,041$ en cinq ans, vers un conforme d'environ $6\,925$ \\
Ordre optimal & sur les $120$ ordres, P1, P4, P2, P3, P5, l'ordre ROOT, distinct du glouton ;
  pire contre optimal $21\,\%$ \\
Ensemble admissible & P1 premier dans $43{,}6\,\%$, P4 dans $37{,}7\,\%$, écart de regret maximal
  $1{,}5\,\%$ ; P5 jamais premier \\
Élasticités & $\xi$ $3{,}79$ ; $\sigma$ $0{,}97$ ; $p_u$ $0{,}82$ ; $\lambda$ $0{,}45$ ; $\varphi$
  $-0{,}44$ ; $g$ $0{,}37$ ; seuil seul $0{,}03$ ; ablation : propagation $-20\,\%$, queue
  $-76\,\%$ \\
Copule & gaussienne à marges égales : $+0{,}1\,\%$ (et $+4{,}8\,\%$ en moyenne sur quatre
  graines) ; Student $+17\,\%$ ; surcoût $6\,760$ contre $10\,927$ ($-38\,\%$) ; interaction
  moyenne $0$ contre $+560$ ; part de l'effet : Formule Standard $0\,\%$, copule $62\,\%$ \\
\end{fichetab}
\end{parcoeur}

\begin{question}{\og Vous dites ne pas connaître la direction, et vous publiez un facteur $3{,}34$. Laquelle est votre réponse ?\fg{}}
\rep{Les deux, car ils ne mesurent pas la même chose. Le facteur $3{,}34$ mesure l'effet des quatre canaux, avec une direction fixée : celle du classeur. La bande mesure ce que coûte l'ignorance de la direction. Le chapitre fait aussi le lien entre son état conforme d'environ $5\,900$ et le socle de $5\,275$ du chapitre~9. Ce sont deux points du même modèle, pas deux mesures du même objet.}
\appui{\S~9.1, folios 60 à 62, encadré ; chapitre~8, folios 56 à 59.}
\end{question}

\begin{question}{\og Pourquoi publier la lecture à quatre canaux, celle qui donne le plus gros chiffre ?\fg{}}
\rep{Parce que c'est la seule où l'état conforme est conforme sur tous les canaux du modèle. Dans les autres lectures, l'entité dite conforme propage encore comme une entité défaillante. Ou bien elle garde sa détection ou son accumulation au niveau non conforme. C'est alors une remédiation partielle. L'appeler « conforme » sous-estimerait l'écart. Je publie aussi les trois autres lectures, comme des remédiations partielles.}
\appui{\S~9.1, folios 60 à 62, encadré.}
\end{question}

\begin{question}{\og Vos canaux isolés somment à $9\,138$, l'écart vaut $14\,139$ : où sont les $5\,001$ ?\fg{}}
\rep{Ils sont dans l'interaction. Elle vient d'une composition presque multiplicative : le produit des quatre facteurs isolés vaut $3{,}47$, pour un facteur publié de $3{,}34$. Si on décompose en additions des effets qui se multiplient, on obtient mécaniquement une interaction positive. Pour la gestion, cela veut dire : chiffrer une remédiation par l'effet de \emph{fermeture}, depuis l'état non conforme. Jamais par la somme des effets isolés.}
\appui{\S~9.1, folios 60 à 62, encadré ; supplément S1, \S~S1.4.}
\end{question}

\begin{question}{\og Y a-t-il un double comptage entre vos canaux ?\fg{}}
\rep{Non. Chaque canal agit sur un objet différent du modèle : le nombre d'incidents, la détection, la propagation et l'accumulation par un prestataire. La colonne des effets de fermeture semble compter deux fois, mais c'est une identité. Sa somme vaut $\sum_S|S|\,m(S)$ : un croisé d'ordre $k$ y est compté $k$ fois. Soit $1\times9\,138+2\times5\,345-3\times691+4\times346=19\,141$. La vraie propriété est l'encadrement $9\,138\le14\,139\le19\,141$.}
\appui{supplément S1, encadré.}
\end{question}

\begin{question}{\og Votre modèle est-il super-additif ?\fg{}}
\rep{Il faut distinguer trois énoncés. Dans un sinistre, les coûts des piliers touchés s'additionnent : c'est une hypothèse de construction. En fonction des piliers non conformes, le capital est presque additif ($R^2=0{,}9945$) : c'est une mesure. En fonction des quatre canaux, il est super-additif ($+35\,\%$) : c'est aussi une mesure. Répondre « oui » tout court serait faux.}
\appui{\S~10.7, folio 76.}
\end{question}

\begin{question}{\og Shapley et Euler ne donnent pas le même classement.\fg{}}
\rep{C'est le résultat. Shapley regarde la \emph{source} de la non-conformité. Il concentre le surcoût sur P1, à $34\,\%$. Euler regarde le pilier \emph{touché}. Il répartit le capital presque uniformément, avec P2 et P4 légèrement en tête. La remédiation se décide du côté de la source. Le capital se loge là où les cascades arrivent.}
\appui{\S~9.3, folios 63 et 64, encadré.}
\end{question}

\begin{question}{\og Votre SCR est-il centré ?\fg{}}
\rep{Non. Partout, c'est le quantile $\mathrm{VaR}_{99{,}5\,\%}(L)$ non centré, au niveau de Solvabilité~II. La perte attendue d'une non-conformité n'est couverte par aucune provision, donc je ne la déduis pas. Pour l'ordre de grandeur : centré, le chiffre de tête tomberait à environ $5\,430\to16\,320$, soit un facteur d'environ $3{,}0$ (calcul hors mémoire).}
\appui{\S~9.3, folios 63 et 64.}
\end{question}

\begin{question}{\og Pourquoi une VaR et pas une TVaR ?\fg{}}
\rep{Parce que Solvabilité~II définit le capital comme une VaR à $99{,}5\,\%$ sur un an. Un travail qui se compare au cadre prudentiel doit parler sa langue. J'utilise la TVaR comme diagnostic, pour l'allocation d'Euler et pour la forme de la queue : $\mathrm{TVaR}/\mathrm{VaR}\approx2{,}1$, contre $2{,}5$ en théorie. L'argument décisif contre elle est l'existence. Avec la calibration PRC, $\hat\xi=1{,}03>1$. La moyenne est alors infinie et la mesure n'existe pas.}
\appui{\S~9.3, folios 63 et 64.}
\end{question}

\begin{question}{\og Quel canal une direction des risques doit-elle maîtriser en premier ?\fg{}}
\rep{La fréquence, et le test de résistance inversé le montre. Seule, elle porte le capital à $10\,377$. Aucun autre canal seul n'atteint $8\,000$. Au-delà de $11\,413$, ce que les trois autres canaux atteignent ensemble, aucune cible n'est atteignable sans elle. La lecture inverse désigne le canal dont la maîtrise empêche les scénarios sévères.}
\appui{\S~9.5, folio 65 ; supplément S1, \S~S1.4.}
\end{question}

\begin{question}{\og Qu'est-ce qu'un dirigeant fait de votre travail lundi matin ?\fg{}}
\rep{Il peut faire trois choses. D'abord, il sait quel pilier remédier en premier. L'ordre optimal n'est pas l'ordre glouton, car les interactions obligent à raisonner sur tout l'horizon : cela représente $21\,\%$ de capital intégré entre le pire ordre et l'ordre optimal. Ensuite, il peut comparer remédier, détenir et transférer sur un seul axe, en euros par an. Enfin, il sait quelle donnée demander à son registre d'incidents.}
\appui{\S~9.5, folio 65 ; \S~9.13, folio 70 ; \S~10.8, folio 76.}
\end{question}

\begin{question}{\og Votre recommandation dépend-elle de paramètres que personne ne sait calibrer ?\fg{}}
\rep{Non. L'ordre de remédiation optimal ne dépend que des SCR par configuration. Il ne change pas avec les taux de transition de la chaîne de Markov, qui sont pourtant les paramètres les plus difficiles à calibrer. Ces taux fixent le calendrier, jamais la séquence.}
\appui{\S~9.6, folio 65, encadré.}
\end{question}

\begin{question}{\og P1 est-il vraiment le premier pilier à remédier ?\fg{}}
\rep{Avec la matrice fixée, Shapley donne $34\,\%$, et l'ordre optimal est P1 puis P4. Dans les directions examinées, P1 et P4 sont premiers dans $43{,}6\,\%$ et $37{,}7\,\%$ des cas, avec un écart de regret de $1{,}5\,\%$. Ils arrivent le plus souvent en tête, mais P2 ou P3 peuvent passer devant ; P5 n'est jamais premier dans cet échantillon. Avec le biais de narration, la priorité passe à P4 à mi-chemin. Une part Shapley, un gain de remédiation et un ordre optimal ne se confondent pas.}
\appui{\S~9.6, folio 65, encadré.}
\end{question}

\begin{question}{\og Votre modèle est-il sensible aux probabilités de propagation, qui ne sont pas observées ?\fg{}}
\rep{C'est le levier le moins sensible. L'élasticité est de $0{,}37$, contre $3{,}79$ pour l'indice de queue. Retirer la propagation fait baisser la VaR de $20\,\%$. Retirer la queue lourde la fait baisser de $76\,\%$. Pour la direction, je ne la stresse pas, je la borne : $1\,839$ en cas d'ignorance totale, soit $34{,}9\,\%$ du socle. La vigilance doit donc porter d'abord sur les valeurs extrêmes.}
\appui{\S~9.7, folios 65 et 66.}
\end{question}

\begin{question}{\og Pourquoi une cascade et pas une copule ?\fg{}}
\rep{Le comparateur à copule gaussienne figée donne un niveau proche de la cascade, à $+0{,}1\,\%$ près, mais un écart entre états de $6\,760$ contre $10\,927$~M\euro. Son interaction moyenne est nulle par construction, contre $+560$ pour la cascade. Ces différences caractérisent le comparateur testé, dont la matrice de corrélation reste fixe entre états. Elles ne démontrent pas qu'une copule dépendant de l'état serait incapable de représenter ces effets.}
\appui{\S~9.8, folio 66.}
\end{question}

\begin{parcoeur}[À savoir par c\oe ur : l'entité, l'agrégation, le transfert]
\begin{fichetab}
Descente & $\lambda(A)=\lambda_c(A/A_c)^{b_\lambda}$, $b_\lambda=0{,}074$ ;
  $m(A)=(A/A_m)^b$, $b=0{,}087$, $m=0{,}854$ à la taille de l'entité notionnelle \\
Échelles & secteur $\lambda=21{,}56$ : $8\,123$ M\euro{} (moyenne $983$) ; seau d'au moins dix
  événements $0{,}210$ : $488$ ; seau toutes firmes $0{,}100$ : $229$ ; $\lambda$ lu à la taille
  $0{,}092$ : $198$ ; les deux canaux composés : $169$ (moyenne $3{,}6$), bande
  $[131{,}5\,;183{,}3]$ \\
Non-linéarité & diviser $\lambda$ par $235$ divise la moyenne par $270$ et le capital par $48$ :
  le quantile est porté par le plus grand sinistre \\
Deux incertitudes & à l'entité, la bande de direction mesure $51{,}7$ M\euro{} ; l'IC de $\xi$
  donne $140{,}9$ à $191{,}6$, soit $51{,}5$ : du même ordre \\
Quatre SFCR 2024 & non-vie A : actif $2\,319$, SCR $425$, besoin $112$ ($26{,}4\,\%$) ; non-vie
  B : $5\,305$, $812$, $136$ ($16{,}7\,\%$) ; vie C : $37\,845$, $1\,587$, $196$ ($12{,}3\,\%$) ;
  vie D : $309\,800$, $14\,800$, $294$ ($2{,}0\,\%$) \\
Extrapolation & un facteur $134$ sur l'actif ne divise la sévérité que par $1{,}53$ et la
  fréquence par $1{,}44$. Le repère conventionnel de $10\,\%$ du SCR est franchi entre $37\,845$ et
  $309\,800$ ; il ne définit pas une validité statistique. Pour l'entité notionnelle ($19\,231$),
  le levier médian ($10{,}0\,\%$) donne $9\,\%$ du SCR. Le besoin $169$ n'est pas un majorant établi \\
Forfait de $3\,\%$ & $54{,}1$ contre $59$ publié ($-8\,\%$), $62{,}1$ contre $39{,}7$ ($+56\,\%$) :
  un ordre de grandeur, ni la charge ni une borne \\
Agrégation & $\tau(\rho,d)=\frac{\sqrt{1+2\rho d+d^2}-1}{d}=\rho+\frac d2(1-\rho^2)+O(d^2)$ ;
  à $\rho=0{,}25$, rapport additif sur corrélé de $3{,}4$ à $3{,}9$ ; l'écart pèse de $0{,}58$ à
  $8{,}21\,\%$ du SCR publié \\
Transfert & coût du capital $6\,\%$, $4{,}75\,\%$ après la directive (UE) 2025/2 ; sinistres sur
  primes $48{,}5\,\%$ ; multiple $\mathrm{SCR}/\mathbb E[L]$ de $47$ à l'entité contre $8{,}3$ au
  secteur ; détention $8{,}0$ M\euro/an ($10{,}1$ à $6\,\%$), $2{,}2$ à $2{,}8$ fois la perte
  attendue ; portée optimale $169$, coût $8{,}0\to2{,}9$ ($-64\,\%$) ; seuil d'inefficacité
  $64\,\%$ ($81{,}5\,\%$ pertes conservées comprises) ; fréquence divisée par deux, valeur posée : coût optimal $-60\,\%$ \\
\end{fichetab}
\end{parcoeur}

\begin{question}{\og Ce montant sectoriel, que vaut-il pour mon entité ?\fg{}}
\rep{La lecture du registre donne $169$~M\euro{}, dans une plage de $131{,}5$ à $183{,}3$ aux sommets. C'est une extrapolation, sans garantie de majoration : la taille adapte l'échelle de la sévérité, pas sa forme. Le calcul utilise la cascade indépendante sans choc commun P4, tandis que le chiffre de tête à quatre canaux emploie la marche avec ce choc. La version Institut ajoute des lectures d'exposition ; elles restent distinctes du résultat du rapport déposé.}
\appui{\S~9.10, folios 67 et 68 ; \S~9.11, folios 68 et 69, encadré.}
\end{question}

\begin{question}{\og Quelle est la limite de votre extrapolation à une entité ?\fg{}}
\rep{La sévérité garde la forme estimée sur de grandes institutions, et aucun coefficient métier n'est estimé dans le moteur. Le rapport au SCR est un repère de taille ; le seuil de $10\,\%$ est conventionnel. Il ne permet ni de déclarer un domaine de validité statistique ni de prouver que le besoin est surestimé. La réserve appelle une donnée d'exposition et une validation à l'échelle de l'entité.}
\appui{\S~9.11, folios 68 et 69.}
\end{question}

\begin{question}{\og Une PME n'a pas un facteur de descente plus fort que le Crédit Agricole. Votre méthode le prévoit-elle ?\fg{}}
\rep{Non, et c'est la limite principale de la descente d'échelle. Les élasticités à la taille sont faibles, $0{,}074$ et $0{,}087$. Une petite entité subit donc presque la sévérité d'un grand groupe, d'où une charge de $26{,}4\,\%$ du SCR pour le plus petit assureur contre $2{,}0\,\%$ pour le plus grand. La méthode ne tient pas compte du métier non plus. Un essai exploratoire, hors rapport, montre qu'à taille égale l'assurance a environ la moitié de la fréquence de la banque et environ $2{,}4$ fois sa sévérité. Pour corriger, il faudrait une échelle selon la taille et le type d'activité, avec un plafond par sinistre lié à l'exposition. Ce plafond serait posé, faute de données sur les petites entités. C'est une piste d'amélioration, pas un résultat.}
\appui{\S~9.10 et \S~9.11, folios 67 à 69 ; essai exploratoire, hors rapport.}
\end{question}

\begin{question}{\og Pourquoi anonymiser des entités dont les rapports sont publics ?\fg{}}
\rep{Par prudence, pas par confidentialité. Le modèle n'utilise que le bilan. L'état de conformité y est \emph{supposé}. Nommer une entité reviendrait à lui attribuer une non-conformité que rien n'observe. La table de correspondance existe dans un script. Le travail reste donc vérifiable par qui a le dépôt.}
\appui{\S~9.11, folios 68 et 69, encadré.}
\end{question}

\begin{question}{\og Pourquoi refusez-vous d'appeler cela un SCR ?\fg{}}
\rep{Parce qu'aucun module DORA n'existe dans la Formule Standard. Ce que je calcule est un besoin de capital ORSA, donc de pilier~2. Le rapport au SCR publié est une mise à l'échelle, pas une part. Diviser un besoin de pilier~2 par le SCR réglementaire n'en fait pas une composante. Écrire « le SCR DORA de telle entité » ferait croire à un objet qui n'existe pas.}
\appui{\S~9.11, folios 68 et 69, encadré ; \S~1.6, folios 18 et 19.}
\end{question}

\begin{question}{\og La Formule Standard n'est-elle pas au moins une borne ?\fg{}}
\rep{Non. Sa réponse à la conformité est nulle par construction : $0\,\%$ de l'effet, contre $62\,\%$ pour le comparateur à copule gaussienne figée et $100\,\%$ pour la cascade. Le repère de $3\,\%$ des provisions s'écarte des modules publiés de $-8\,\%$ à $+56\,\%$. Je ne le cite donc que comme ordre de grandeur. La coïncidence de $488$ contre $450$ au seuil de dix événements ne résiste pas à la correction d'échelle.}
\appui{\S~9.9, folios 66 et 67 ; \S~9.11, folios 68 et 69.}
\end{question}

\begin{question}{\og Que retient le capital total de votre écart ?\fg{}}
\rep{Tout, avec la convention de Solvabilité~II. La charge opérationnelle s'ajoute au BSCR sans passer par la matrice de corrélation. Si l'écart était agrégé comme un module corrélé, il ne compterait qu'à hauteur de $\tau$, qui tend vers $\rho$ pour un petit écart. À $\rho=0{,}25$, l'écart de traitement vaut de $3{,}4$ à $3{,}9$. Un euro économisé sur DORA vaut donc trois à quatre fois plus, en capital, qu'un euro sur un risque qui se diversifie.}
\appui{\S~9.12, folios 69 et 70, encadré.}
\end{question}

\begin{question}{\og Détenir ou transférer ?\fg{}}
\rep{À l'échelle d'une entité, détenir le capital coûte $8{,}0$~M\euro{} par an au taux révisé. C'est de $2{,}2$ à $2{,}8$ fois la perte attendue : immobiliser coûte plus que le risque lui-même. Avec un traité en excès de perte, de portée égale au besoin du modèle ($169$), le coût de financement tombe à $2{,}9$, soit $-64\,\%$. En comptant les pertes conservées, il passe de $11{,}67$ à $5{,}12$, soit $-56\,\%$. C'est une borne supérieure. La capacité du marché, les exclusions et le risque de contrepartie la réduisent. Face au meilleur traité, le transfert reste préférable tant que moins de $64\,\%$ de la portée achetée reste sans allègement de capital ($81{,}5\,\%$ pertes conservées comprises). Enfin, une entité conforme (fréquence de $0{,}092$ à $0{,}046$) voit ce coût optimal baisser de $60\,\%$.}
\appui{\S~9.13, folio 70.}
\end{question}

\begin{question}{\og Votre bande d'identification est-elle informative, ou trop large pour rien apprendre ?\fg{}}
\rep{Elle est informative, et la comparaison le montre. À l'échelle d'une entité, elle mesure $51{,}7$~M\euro. Faire varier seulement l'indice de queue sur son intervalle à $90\,\%$ déplace le capital de $51{,}5$. Ignorer la direction coûte donc autant qu'ignorer la queue. Une meilleure spécification de reporting a le même poids qu'une meilleure estimation de la queue.}
\appui{\S~9.10, folios 67 et 68, encadré.}
\end{question}

\lepiege{Le multiple de capital de $47$ à l'entité contre $8{,}3$ au secteur n'est pas une
incohérence : c'est la mutualisation lue à l'envers, le capital d'une entité isolée étant dominé
par un sinistre unique, celui du secteur par l'accumulation.}

% ---------------------------------------------------------------------
\section*{Chapitre 10. Robustesse, limites et recommandations (sections 10.1 à 10.7)}
\addcontentsline{toc}{section}{Chapitre 10. Robustesse, limites et recommandations (sections 10.1 à 10.7)}
\reperes{folios 71 à 77. 10.1 la table de robustesse (71), 10.2 limites du domaine et de ce travail (72), 10.3 trois étages (73), 10.4 le moteur contrôlé (74), 10.5 VaR prédictive et robuste (74), 10.6 le défaut de calibration (75), 10.7 défaillances simultanées et additivité des coûts (76).}

\enbref{La robustesse en une table, la séparation entre ce que la donnée interdit à quiconque et ce
que ce travail a choisi, les trois étages d'incertitude sur une même échelle, un moteur contrôlé
par trois chemins, le choix du chiffre reporté et le registre des écarts chiffrés plutôt que
corrigés.}

\begin{parcoeur}[À savoir par c\oe ur : incertitude et moteur]
\begin{fichetab}
Verdicts & robuste : l'effet DORA, la priorité P1 et P4, le rôle de l'asymétrie, l'absence de
  Hawkes, les lois du socle hors échantillon ; voulu : le capital réagit à la conformité
  ($\times3{,}34$) ; assumé : la distribution large, le niveau de sévérité hors échantillon,
  l'additivité des coûts \\
Deux natures & limites du domaine (VERIS $0{,}12\,\%$, $1{,}19\,\%$, un incident ; coût
  multi-piliers inobservé ; sévérité en euros seulement dans les bases bancaires) contre limites
  de ce travail (élicitation sans réponse, gel, source citée, deux écarts publiés) \\
Même échelle & charge annuelle : paramètre $[3\,042\,;27\,152]$, $\times8{,}93$ ; identification
  $[6\,858\,;8\,697]$, $\times1{,}27$ ; modèle, dépendance $[5\,322\,;9\,806]$, $\times1{,}84$ ;
  modèle, famille de queue $[4\,330\,;39\,901]$, $\times9{,}21$ \\
Bande reportée & paramètre plus identification ; le modèle en sort, publié comme axe de
  sensibilité, parce qu'un choix de famille ne se moyenne pas \\
Famille & lognormale $4\,330$ contre GPD à $\xi=0{,}90$ $24\,016$ ($\times5{,}5$) ; cascade
  $8\,318$, entre gaussienne et Gumbel ; plafond de $40$ M\euro{} : $495$ \\
Hawkes & une cascade sans auto-excitation donne $n=0{,}480$ au jour contre $0{,}551$ mesuré ;
  pic à $3{,}3$~h ; $0{,}551$ tombe vers zéro si les co-occurrences du jour passent en exogène \\
Moteur & Monte-Carlo $6\,049$ et $20\,188$ ; Fourier $5\,945$ et $19\,890$ ($-1{,}7$ et
  $-1{,}5\,\%$, soit $0{,}75$ et $0{,}59$ erreur type) ; perte unique au second ordre $-9{,}9$ et
  $-7{,}7\,\%$, au premier $-17{,}5$ et $-26{,}0\,\%$ \\
\end{fichetab}
\end{parcoeur}

\begin{parcoeur}[À savoir par c\oe ur : chiffre reporté et écarts]
\begin{fichetab}
Six postures & point $657$ ; prédictive $645\pm7$ ($648$ ailleurs, le même nombre) ; robuste à
  $\beta=90\,\%$ $939\pm13$, à $95\,\%$ $1\,037\pm9$, à $99\,\%$ $1\,247\pm24$ ; pire cas de
  modèle ($\xi=0{,}90$) $1\,331$ \\
Robuste à $95\,\%$ & c'est la borne haute de l'IC à $90\,\%$ déjà publié, $1\,037$ : la même
  quantité sous deux noms ; la prédictive ne s'écarte du point qu'en profondeur ($+4\,\%$ à
  $99{,}9\,\%$) \\
Posture retenue & le point avec sa bande, pour trois raisons : réglementaire, de mesure, de
  précision \\
$p_u$ gelé & $86/570\simeq0{,}1509$ sur 2000--2026, contre $91/582\simeq0{,}1564$ sur l'échantillon complet ; quantile $662{,}78\to678{,}60$
  si seul le taux change ($+2{,}4\,\%$, $15{,}8$ M\euro) \\
Deux colonnes & limites de calibration : périmètres différents pour $p_u$, couverture $86\,\%$, dérive
  (au moins $+24{,}7\,\%$) ; choix posés : $a=0{,}60$ contre $0{,}22$,
  $\varphi=9{,}20$ contre $2{,}24$ ; périmètre de sens opposés : attritionnel absent, sévérité non
  plafonnée \\
Multi-piliers & $62{,}31\,\%$ des sinistres à l'état non conforme contre $31{,}15\,\%$ à l'état
  conforme, pour la cascade seule ; $73{,}25\,\%$ à l'état non conforme avec le choc P4.
  Moyennes : $1{,}931$ contre $1{,}380$ pour la cascade seule ; $2{,}44$ avec le choc P4 \\
Additivité des coûts & perte $(\sum_jX_j)\,k^{\theta-1}$ ; pour $\theta\in[0{,}70\,;1{,}30]$,
  écart de $9\,706$ à $20\,487$ ($76\,\%$ de l'écart, quinze fois l'écart-type de $734$ entre graines) ;
  élasticités $0{,}95$ (NC) contre $0{,}39$ (C), facteur $2{,}42$ ; sens indéterminé \\
Plafond et saturation & le plafond amortit, l'interaction change de signe à $\theta=1$ ; à
  $\kappa=0{,}5\,\%$, $96{,}2$ M\euro{} pour tous les $\theta$ ; bande conjointe $177{,}8$ sous la
  somme des deux, $203{,}7$ \\
\end{fichetab}
\end{parcoeur}

\begin{question}{\og Vos limites ne sont-elles pas des insuffisances de votre travail ?\fg{}}
\rep{Il y a deux sortes de limites. Les unes tiennent aux sources examinées, qui ne suffisent pas à identifier la direction. Dans l'extrait VERIS étudié, le champ du domaine de contrôle n'est rempli que dans $0{,}12\,\%$ des cas. Les autres sont des choix de ce travail : le gel, la source citée, l'élicitation restée sans réponse, deux écarts publiés plutôt que corrigés. Les premières demandent des observations plus détaillées sur les piliers touchés et l'ordre de propagation. Les secondes demandent de reprendre l'analyse.}
\appui{\S~10.2, folios 72 et 73.}
\end{question}

\begin{question}{\og Des trois étages d'incertitude, lequel domine ?\fg{}}
\rep{Sur la charge annuelle, le paramètre et le modèle pèsent à peu près autant : facteurs $8{,}93$ et $9{,}21$. Tous deux dépassent largement l'identification, à $1{,}27$. Le facteur $2{,}5$ souvent cité porte sur le quantile d'un seul sinistre. Le comparer tel quel à l'étage de modèle, ce serait comparer deux échelles différentes. La bande reportée cumule le paramètre et l'identification. L'étage de modèle est publié à part, comme un axe de sensibilité.}
\appui{\S~10.3, folios 73 et 74.}
\end{question}

\begin{question}{\og Pourquoi sortir l'incertitude de modèle de la bande ?\fg{}}
\rep{Parce que les deux premiers étages sont statistiques : un intervalle peut les porter. Le troisième est d'une autre nature, il tient au choix de la famille de lois, et un choix ne se moyenne pas. Sur l'axe de la queue, il place $\xi$ à une valeur posée, hors de l'intervalle bootstrap. Le cumuler mélangerait un scénario et une borne de confiance dans le même objet. Je le publie quand même deux fois : comme pire cas et comme bande de dépendance.}
\appui{\S~10.3, folios 73 et 74.}
\end{question}

\begin{question}{\og Et si vous aviez pris une autre méthode ?\fg{}}
\rep{Chaque brique a des alternatives : fréquence, sévérité, dépendance, mesure, agrégation. Je ne prétends pas avoir la seule bonne méthode. J'ai testé les concurrents les plus sérieux et j'ai borné l'effet de ce choix. Le niveau bouge, mais il reste dans la bande. L'ordre d'amorce des piliers (P1, P4, P2, P3, P5) ne change pas d'une famille à l'autre.}
\appui{\S~10.3, folios 73 et 74, encadré.}
\end{question}

\begin{question}{\og Pourquoi pas un processus de Hawkes ?\fg{}}
\rep{Je l'ai testé sur des données datées, contre les noyaux de la littérature. L'excitation se fait dans la journée, avec un pic à $3{,}3$~heures. Le ratio de branchement de $0{,}551$ tombe presque à zéro dès que les co-occurrences du même jour sont traitées comme exogènes. Surtout, une cascade sans aucune auto-excitation donne déjà $0{,}480$ au jour. À la résolution disponible, on ne peut pas distinguer les deux. Je garde la cascade parce qu'elle est plus simple, et qu'elle porte en plus la structure entre piliers. C'est une équivalence observationnelle, pas une réfutation.}
\appui{\S~10.1, folios 71 et 72 ; \S~10.3, folios 73 et 74.}
\end{question}

\begin{question}{\og Comment savez-vous que votre moteur de simulation ne se trompe pas ?\fg{}}
\rep{Je le contrôle par des chemins qui ne partagent ni le code ni l'aléa. L'inversion de Fourier de la loi composée s'écarte de $-1{,}7$ et $-1{,}5\,\%$. C'est moins d'une erreur type, donc l'écart n'est pas significatif. L'approximation par perte unique reste à moins de $10\,\%$ au second ordre. Limite : ces chemins partagent les mêmes lois et les mêmes tables de cascade. Ils vérifient l'agrégation, pas la vérité du modèle.}
\appui{\S~10.4, folio 74.}
\end{question}

\begin{question}{\og Quel nombre reportez-vous : le point, le prédictif ou le robuste ?\fg{}}
\rep{Le point, avec sa bande. J'ai trois raisons. La première est réglementaire : une borne haute sur un ensemble d'ambiguïté est un supremum sur une famille de lois, pas un quantile de la loi de perte. La deuxième est la mesure : la prédictive vaut $645$ contre $657$, soit deux fois le bruit. La troisième est la précision : les postures robustes sont les moins reproductibles. De plus, la robuste à $95\,\%$ est déjà publiée : c'est la borne haute de l'intervalle, $1\,037$.}
\appui{\S~10.5, folio 74.}
\end{question}

\begin{question}{\og \og Le quantile d'un objet incertain est mauvais\fg{} : que répondez-vous ?\fg{}}
\rep{Que la remarque est juste, mais pas pour la raison attendue. Intégrer l'incertitude du paramètre avant de prendre le quantile ne déplace pas le point à $99{,}5\,\%$ ($645$ contre $657$). L'écart n'apparaît que plus loin dans la queue : $+4\,\%$ à $99{,}9\,\%$. La fragilité est dans la largeur de la bande, pas dans le point. La règle est donc : ne jamais donner le point seul, toujours avec sa bande.}
\appui{\S~10.5, folio 74.}
\end{question}

\begin{question}{\og Vous savez que votre taux de dépassement est incohérent avec votre seuil, et vous ne corrigez pas ?\fg{}}
\rep{Le taux de $0{,}1509$ correspond à $86$ excès sur $570$ pertes entre 2000 et 2026, au seuil de $20{,}03$ millions. L'échantillon complet ajoute douze pertes antérieures à 2000 : il compte $91$ excès sur $582$ pertes, soit $0{,}1564$. La description du périmètre de calibration était incomplète. Si seul le taux change, le quantile de sévérité passe de $662{,}78$ à $678{,}60$ millions, soit $+2{,}4\,\%$. C'est une sensibilité à paramètres fixés. Une recalibration complète exige aussi de réajuster la forme et l'échelle ; son effet sur le capital annuel et le rapport entre états doit être calculé.}
\appui{\S~10.6, folio 75.}
\end{question}

\begin{question}{\og Vous avez gelé votre calibration en cours de route. N'est-ce pas un aveu ?\fg{}}
\rep{C'est une décision de conduite de projet, datée. Quand on a publié beaucoup de résultats, un modèle qui bouge encore rend faux, sans le dire, ce qui est déjà écrit. Le gel sépare deux choses. La correction d'erreur reste autorisée et je la pratique. La recalibration ne l'est plus. Le prix du gel est publié : les écarts connus sont chiffrés au lieu d'être corrigés.}
\appui{\S~10.6, folio 75 ; enseignements du stage, folios 80 à 82.}
\end{question}

\begin{question}{\og Vos erreurs vont-elles toutes dans le même sens ?\fg{}}
\rep{Il faut distinguer les sensibilités à paramètres fixés et une recalibration complète. Relever seulement le taux de dépassement augmente le quantile de sévérité. Mais le taux gelé correspond à un autre périmètre, et une recalibration complète réajusterait aussi la forme et l'échelle. Je ne peux donc pas conclure que tous les écarts sous-estiment le capital. La couverture de l'intervalle concerne l'incertitude ; la dérive de l'échelle, la charge systémique et la surdispersion sont des diagnostics distincts.}
\appui{\S~10.6, folio 75.}
\end{question}

\begin{question}{\og Les défaillances simultanées de plusieurs piliers ne sont pas traitées.\fg{}}
\rep{La prémisse est fausse : c'est une sortie du modèle. Pour la cascade seule, hors choc P4, $62{,}31\,\%$ des sinistres touchent plus d'un pilier à l'état non conforme, contre $31{,}15\,\%$ à l'état conforme. Avec le choc P4, cette part atteint $73{,}25\,\%$ à l'état non conforme. Ces probabilités sont calculées par énumération, sans bruit Monte-Carlo. Ce qui reste une hypothèse, c'est la façon dont les coûts de ces piliers s'additionnent.}
\appui{\S~10.7, folio 76 ; diagnostic 118 pour la variante avec choc P4.}
\end{question}

\begin{question}{\og Quelle est l'hypothèse la plus lourde de votre modèle ?\fg{}}
\rep{L'additivité des coûts dans un sinistre qui touche plusieurs piliers. Si on la relâche avec un exposant, l'écart entre états passe de $9\,706$ à $20\,487$~M\euro, soit $76\,\%$ de l'écart publié. Elle pèse plus à l'état non conforme : élasticité $0{,}95$ contre $0{,}39$. On ne peut pas dire dans quel sens elle joue. La mutualisation de la remédiation et la saturation de la capacité tirent en sens contraires. Ce qui tient : le signe et l'ordre de grandeur de l'écart, sur toute la plage.}
\appui{\S~10.7, folio 76.}
\end{question}

% ---------------------------------------------------------------------
\section*{Section 10.8. Recommandations sur la donnée manquante}
\addcontentsline{toc}{section}{Section 10.8. Recommandations sur la donnée manquante}
\reperes{folios 76 et 77. 10.8.1 le registre, 10.8.2 ce que l'information vaut, 10.8.3 quelle dépendance documenter.}

\enbref{La limite d'identification se retourne en spécification de reporting, chiffrée par la
valeur d'information et hiérarchisée par dépendance : c'est le seul endroit du mémoire où une
limite produit un livrable.}

\begin{parcoeur}
\begin{fichetab}
Le registre & horodater la survenance au mois ou à la semaine ; catégoriser par domaine de
  contrôle ; consolider une cause commune en un événement à $N$ victimes ; consigner tous les
  piliers qui ont cédé ($62{,}31\,\%$ de multi-piliers sous $W$, $50{,}77\,\%$ sous $W^\top$) \\
Instruments & le registre d'information du règlement d'exécution 2024/2956 sert la troisième
  exigence ; la notification des incidents majeurs porte les trois autres, par le format de
  l'horodatage et la taxonomie \\
VERIS & le champ du contrôle défaillant et la chronologie en quatre étapes existent : rendre
  obligatoires deux champs normalisés, pour un coût de conception nul \\
Valeur & $L(\mathcal V)=\sup_{\mathcal V}\mathrm{SCR}-\inf_{\mathcal V}\mathrm{SCR}$ ;
  $V(\mathcal I)=1-L(\mathcal W_{\mathcal I})/L(\mathcal W)\in[0,1]$ ; monotone et sans
  dimension \\
$32{,}7\,\%$ & réduction moyenne de la plage ($22$ à $44\,\%$ selon la graine, quatre graines) si l'on
  retenait le sens que sept rapports documentent sur six dépendances, alors que deux seulement sont
  tranchées : un indicateur, pas un acquis ; grille commune, chaque coefficient à $-S$, $0$ ou $S$ \\
$11{,}2$ à $13{,}3\,\%$ & plage levée en moyenne par chacune des quatre dépendances de tête (dont
  gouvernance-incidents, $13{,}3$) ; les cinq dernières, de $0{,}7$ à $2{,}3\,\%$ ; ordre variable
  entre graines \\
Amplitude & même les dix sens connus ($57{,}9\,\%$), l'amplitude resterait inconnue : un sens connu
  laisse le coefficient libre entre zéro et son maximum ; l'identification partielle est
  structurelle \\
\end{fichetab}
\end{parcoeur}

\begin{question}{\og Que répondez-vous à un superviseur ?\fg{}}
\rep{Je propose une recommandation chiffrée, à laquelle un superviseur peut s'opposer. Dans la notification des incidents majeurs sous DORA, rendre obligatoires deux champs \emph{déjà normalisés} dans VERIS : le domaine de contrôle défaillant et la chronologie. Le coût de conception est nul. Si on ne pouvait exiger qu'une seule dépendance, la donnée ne dit pas laquelle. Quatre dépendances, dont celle entre gouvernance et gestion des incidents, lèvent chacune de $11$ à $13\,\%$ de la plage en moyenne. Leur ordre change selon les graines.}
\appui{\S~10.8, folio 76.}
\end{question}

\begin{question}{\og $32{,}7\,\%$ de resserrement : c'est un résultat ?\fg{}}
\rep{Non, ce n'est pas un résultat. C'est la valeur moyenne d'une information de cette étendue. Elle va de $22$ à $44\,\%$ selon la graine. Les sept rapports donnent le sens de six dépendances, mais je n'en revendique que deux, paire par paire. Si on retenait ces six sens, la largeur baisserait de $32{,}7\,\%$ en moyenne. C'est l'ordre de grandeur de ce qu'un registre bien conçu apporterait, pas un acquis. Une première version annonçait $66\,\%$. Elle fixait les amplitudes à leur maximum, alors qu'un sens connu les laisse libres. Je l'ai refaite sur une grille commune.}
\appui{\S~10.8.2, folio 76.}
\end{question}

\begin{question}{\og DORA ne prévoit-il pas déjà ce registre ?\fg{}}
\rep{En partie. Le registre d'information contient les accords contractuels avec les prestataires, donc la dépendance \emph{déclarée}. Il sert à repérer les causes communes. Il ne donne pas l'ordre des défaillances. La notification des incidents majeurs existe déjà. Il s'agit d'en fixer le format d'horodatage et la taxonomie. On ne crée pas une nouvelle collecte.}
\appui{\S~10.8.1, folio 76.}
\end{question}

\begin{question}{\og Si l'on connaissait tous les sens, auriez-vous un point ?\fg{}}
\rep{Non. Connaître le sens d'un lien ne donne pas sa force. Même avec les dix directions établies ($57{,}9\,\%$ de la plage levée), on ne connaîtrait pas l'amplitude. L'identification partielle n'est pas un état provisoire que les données finiraient par lever. C'est une raison de plus de ne pas promettre à un superviseur plus qu'un registre ne peut tenir.}
\appui{\S~10.8.2, folio 76, encadré ; \S~7.10, folios 50 à 55.}
\end{question}

\begin{question}{\og Votre registre est-il réaliste ?\fg{}}
\rep{Les praticiens interrogés le jugent réaliste. Ils demandent des règles de classement précises, car le classement par domaine peut être subjectif. Ils signalent aussi un point de vigilance sur la confidentialité. L'un d'eux décrit un usage déjà en place : un tableau de bord des incidents par prestataire, sur douze mois, dans une analyse de risque au titre de DORA.}
\appui{supplément S4.}
\end{question}

% ---------------------------------------------------------------------
\section*{Conclusion}
\addcontentsline{toc}{section}{Conclusion}
\reperes{folios 78 et 79. Résultats établis, non-identifiabilité démontrée de la direction,
conséquences pour le régulateur, conséquences pour l'entité, limites, ouvertures.}

\enbref{Ce qui est établi (le chiffre, le mécanisme, un acquis de méthode), ce qui est démontré non
mesurable (la direction, convertie en bornes), ce qui en découle pour le régulateur et pour
l'entité, puis les limites et trois ouvertures.}

\begin{parcoeur}
\begin{fichetab}
Fréquence & seule, $10\,377$, que la forme fermée de la perte unique retrouve à moins de $1\,\%$ \\
Méthode & le dispositif de vérification a trouvé trois valeurs fausses qu'aucune relecture
  n'avait vues ; il ne vérifie ni qu'un nombre est au bon endroit, ni ce que la phrase en fait \\
Détenir, transférer & rapport du coût de détention au prix du transfert : $2{,}78$ au taux
  révisé, $3{,}51$ au taux historique ; ses deux termes sont proportionnels au même SCR ; l'effet
  de forme est borné à $+36\,\%$ sur une plage de SCR d'un facteur quatre \\
Limites & élicitation sans réponse ; niveaux de propagation posés ; additivité des coûts
  ($0{,}95$ contre $0{,}39$) ; l'ordre entre P1 et P4 dépend de la direction \\
Ouvertures & l'élicitation ; le pilier des tiers sur données clients ; une criticité au sens de
  l'audit, le risque résiduel après contrôles \\
\end{fichetab}
\end{parcoeur}

\begin{question}{\og Qu'apporte votre travail à la science actuarielle ?\fg{}}
\rep{Un mécanisme de dépendance dirigée, dont trois propriétés sont démontrées. L'une d'elles : une matrice et sa transposée, que les statistiques disponibles ne distinguent pas, donnent un capital et une décision différents. Aussi une frontière d'identification établie sur des données réelles. Et une méthode venue de l'économétrie, l'identification partielle, appliquée à une exigence de capital.}
\appui{conclusion, folios 78 et 79.}
\end{question}

\begin{question}{\og Quelle est la faiblesse de votre travail ?\fg{}}
\rep{Je la dis sans détour, parce que le jury la connaît et qu'esquiver coûte plus que l'avouer. Le niveau absolu du capital n'est pas une mesure. L'ordre entre les deux piliers de tête dépend d'une direction que je ne sais pas identifier. Ce que je peux dire, c'est que ces deux piliers dominent les trois autres, et combien coûte l'ignorance.}
\appui{conclusion, folios 78 et 79.}
\end{question}

\begin{question}{\og Pourquoi votre rapport entre détention et transfert serait-il plus robuste que le niveau ?\fg{}}
\rep{C'est une identité. Les deux termes sont proportionnels au même SCR, qui se simplifie. Un simple changement d'échelle ne change donc pas le rapport. Il vaut $2{,}78$ au taux révisé et $3{,}51$ au taux historique. Il reste un effet de forme, car la descente d'échelle déplace séparément la fréquence et la sévérité. Cet effet est borné à $+36\,\%$. Le rapport est plus robuste que le niveau, mais pas invariant.}
\appui{conclusion, folios 78 et 79.}
\end{question}

% ---------------------------------------------------------------------
\section*{Enseignements du stage}
\addcontentsline{toc}{section}{Enseignements du stage}
\reperes{folios 80 à 82, dans la seule version ENSAE. Le cadre de la mission, les enseignements de
l'ENSAE, l'encadrement et les tiers, bilan personnel, recul sur la démarche.}

\enbref{Le recul sur la mission, à la première personne et sans aucun nombre : quatre contraintes de
cadre, les enseignements mobilisés, la part de l'encadrement et des tiers, et ce que je ferais
autrement.}

\begin{parcoeur}
\begin{fichetab}
Quatre contraintes & deux livrables, le mémoire et l'analyse LUCY co-écrite avec Hugo Rapior ; une
  donnée sous licence, où la traçabilité remplace la reproductibilité complète ; des entités
  anonymisées ; le gel de la calibration \\
Enseignements & valeurs extrêmes, modèles à latente gaussienne, processus stochastiques (Markov,
  branchement), statistique inférentielle, Monte-Carlo, jeux coopératifs (Shapley), assurance
  non-vie ; les deux qui ont le plus pesé : l'économétrie de l'identification, et $\xi$ comme
  condition d'existence des moments \\
Encadrement & Hugo Rapior, point hebdomadaire en diapositives ; Caroline Hillairet, sur le
  traitement de l'incertitude (quantile prédictif, sensibilité à la propagation) et sur Hawkes \\
Tiers & un second codage sollicité, non reçu ; trois praticiens ont répondu, non nommés \\
Autrement & le dispositif de vérification dès le début ; le gel de la calibration plus tôt \\
\end{fichetab}
\end{parcoeur}

\begin{question}{\og Qu'est-ce que l'ENSAE vous a apporté qui vous a servi ?\fg{}}
\rep{Surtout deux choses. D'abord l'économétrie de l'identification : reconnaître qu'un paramètre n'est pas identifié, et non pas mal estimé, a changé ma méthode. Ensuite, la lecture de l'indice de queue comme condition d'existence des moments. Cela m'a permis d'écarter des variantes par un argument d'existence plutôt que de degré. La programmation scientifique a compté autant : tenir une chaîne de scripts reproductible pendant des mois est un problème d'ingénierie.}
\appui{enseignements du stage, folios 80 à 82.}
\end{question}

\begin{question}{\og Quelle a été la part de votre maître de stage et de votre tutrice ?\fg{}}
\rep{Leurs demandes ne portaient pas sur des calculs, et c'est ce qui les a rendues utiles. Hugo Rapior m'a fait retirer des lectures que la mesure ne soutenait pas. Il m'a aussi fait préciser un vocabulaire qui cachait une borne prise à l'envers. Caroline Hillairet a travaillé sur le traitement de l'incertitude. Cela a donné deux sections : le quantile prédictif et la sensibilité aux probabilités de propagation. J'ai aussi testé le processus de Hawkes contre ses propres noyaux.}
\appui{enseignements du stage, folios 80 à 82.}
\end{question}

\begin{question}{\og Que feriez-vous autrement ?\fg{}}
\rep{Deux choses. Je mettrais le dispositif de vérification en place dès le début, et non à mi-parcours. Il a fallu rattacher chaque section à ses scripts après coup. Et je gèlerais la calibration plus tôt. Pendant des semaines, des recalibrations partielles ont fait bouger des nombres déjà écrits.}
\appui{enseignements du stage, folios 80 à 82.}
\end{question}

\begin{question}{\og Qu'est-ce qui distingue une mission d'un exercice académique ?\fg{}}
\rep{En conseil, un résultat doit pouvoir être contrôlé par quelqu'un qui n'a pas fait le calcul. Je rattache chaque nombre publié à la sortie du script qui le produit. Cela répond à cette exigence, et cela a pris une part importante du temps.}
\appui{enseignements du stage, folios 80 à 82.}
\end{question}

% ---------------------------------------------------------------------
\section*{Les annexes et le supplément}
\addcontentsline{toc}{section}{Les annexes et le supplément}
\reperes{Dans le rapport : annexe A démonstrations du modèle de cascade, folio 87 ; annexe B
paramètres et notations, folio 91 ; annexe C usage de l'IA générative, folio 95. Hors du rapport,
dans le supplément en ligne : S1 adaptations par pilier, S2 pistes pour de futurs mémoires, S3
questionnaire d'élicitation, S4 diagnostics secondaires (dont le protocole d'élicitation et les
relectures de praticiens).}

\subsection*{Annexe A. Les démonstrations à savoir refaire au tableau}

\begin{parcoeur}
\begin{fichetab}
Forme réduite & $\mathbb E\,\mathcal M_n(a,k)=(W^n)_{ak}$ ; la série de Neumann converge vers
  $(I-W)^{-1}$ si et seulement si $\rho(W)<1$ ; $\mathbf 1\{k\in S_a\}\le\sum_n\mathcal M_n(a,k)$
  donne la majoration de $\mathbb E|S_a|$ \\
Sous-criticité & $\rho(W)\le\lVert W\rVert_\infty=\max_j\sum_kW_{jk}\le g<1$ \\
Loi exacte & deux événements sur des arêtes disjointes, donc indépendants : $B(A)$ pour les
  fuites, $R(A)$ par récurrence sur les sous-ensembles \\
Monotonie & couplage par uniformes communes : $U_{jk}<W_{jk}\Rightarrow U_{jk}<W'_{jk}$, donc
  $S_a(W)\subseteq S_a(W')$ presque sûrement \\
Ordre & pour des chaînes de longueur deux, l'invariance équivaut à
  $T_{ij}=T_{ji}$ \\
Co-occurrence & contre-exemple de la fourche : $W_{12}=W_{13}=w$ ; sous $W$, la paire $\{2,3\}$
  co-occurre avec la probabilité $rw^2$ en cascade indépendante, sous $W^\top$ jamais, pour la même partie symétrique \\
Shapley & poids $\frac{s!(n-s-1)!}{n!}$ : moyenne sur les $n!$ ordres ; la somme des
  contributions le long d'un ordre est télescopique, d'où $\sum_i\varphi_i=v(\text{tous})$ \\
\end{fichetab}
\end{parcoeur}

\begin{question}{\og Montrez que $W$ et $W^\top$ peuvent être distingués.\fg{}}
\rep{Prenons trois piliers, avec $W_{12}=W_{13}=w$ et tout le reste à zéro. Avec $W$, une amorce en~1 touche~2 et~3 ensemble avec la probabilité $w^2$ : la paire $\{2,3\}$ co-occurre. Avec $W^\top$, les arêtes vont de~2 et de~3 vers~1, et aucun sinistre ne touche~2 et~3 ensemble. La partie symétrique est la même, mais les lois sont différentes. C'est la structure en V de l'identification des graphes causaux.}
\appui{annexe~A, folio 87.}
\end{question}

\begin{question}{\og Pourquoi les parts de Shapley somment-elles exactement à $100\,\%$ ?\fg{}}
\rep{C'est l'efficience. Le poids de Shapley fait de $\varphi_i$ la moyenne, sur les $n!$ ordres d'arrivée, de la contribution marginale de $i$. Pour un ordre donné, la somme des contributions se simplifie et vaut $v(\text{tous})-v(\emptyset)$. En faisant la moyenne, on obtient $\sum_i\varphi_i=v(\text{tous})$. Limite : c'est un partage exact du surcoût. Ce n'est pas une prévision de ce qu'on gagnerait en retirant un seul joueur.}
\appui{annexe~A, folio 87.}
\end{question}

\subsection*{Supplément S1. Adaptations par pilier}

\begin{parcoeur}
\begin{fichetab}
Queue par pilier & $\xi$ de $0{,}87$ à $1{,}37$ selon la catégorie bâloise ; $111$ observations
  pour la catégorie TIC, sous le seuil de $150$ ; $120$ assignations : capital $+55$ à $+96\,\%$ ;
  pilier de tête dominant dans $83\,\%$ \\
P4 en choc commun & intensité bornée dans $[0\,;0{,}68]$ ; $3{,}7$ piliers emportés d'un coup,
  $+19\,\%$ de capital ; au moins trois piliers touchés : $0{,}22$ en cascade, $0{,}90$ en choc
  commun ; co-occurrence sans direction, mais intensité non identifiée \\
Seuil et copule & $K_4=1{,}03$ pour $p_4=0{,}151$ ; $\lambda_U=2-2^{1/\theta}=0{,}53$ à
  $\theta=1{,}8$ ; cinq piliers avec P4 : $0{,}17$ en gaussien, $0{,}39$ en Gumbel \\
Temps d'arrêt & $\Prob(\tau\le T)=2\,\Phi(-z^\star/\sqrt T)$, par le principe de réflexion
  (Markov fort et symétrie, pas la martingalité) ; $0{,}14$ à $\rho=0{,}50$, $0{,}23$ à
  $\rho=0{,}75$ \\
Ancrages & ESAs : $29\,\%$ des $3\,383$ incidents majeurs de 2025 d'origine tierce ; BCE : plus
  de $30\,\%$ du budget des grandes banques à dix fournisseurs ; $19$ prestataires critiques
  désignés \\
Détection & VCDB ($n=278$) : $47\,\%$ le jour même, $8\,\%$ en semaines, $45\,\%$ en mois ; canal
  seul de $1\,775$ à $8\,052$ ; le multiplicateur $0{,}85$ et $1{,}20$ correspond à
  $\eta\approx0{,}36$ \\
Coin testé & sur les seize configurations et $65$ paires emboîtées, le capital estimé respecte l'ordre, graine
  par graine ; faux trajectoire par trajectoire ($16{,}2\,\%$ et $30{,}7\,\%$) ; un coin peut être
  infaisable \\
Paires & fréquence et détection $1\,739\pm708$, fréquence et accumulation $1\,666\pm354$,
  fréquence et propagation $1\,182\pm444$ ($86\,\%$ de l'ordre~2) ; propagation et accumulation
  $-330\pm242$, substituts \\
Retour & portage $848$ M\euro/an à $6\,\%$ ($672$ à $4{,}75\,\%$) ; sinistralité évitée $3\,247$,
  $3{,}8$ fois plus ; retour de $6$ à $35$ ans, un majorant \\
\end{fichetab}
\end{parcoeur}

\begin{question}{\og Votre retour sur investissement de $6$ à $35$ ans est-il un plancher ?\fg{}}
\rep{C'est un majorant de la durée. Je ne compte que le portage du capital libéré : $848$~M\euro{} par an. La sinistralité évitée, $3\,247$ par an, vaut $3{,}8$ fois plus. Sous-estimer le bénéfice allonge la durée. La durée réelle est donc plus courte, jamais plus longue. En additionnant les deux, on obtiendrait de $1$ à $7$ ans. Mais on mélangerait un flux de résultat et un coût du capital.}
\appui{supplément S1, \S~S1.4.}
\end{question}

\begin{question}{\og Un test de résistance sur vos quatre canaux ?\fg{}}
\rep{Parmi les seize configurations évaluées, le coin non conforme donne le quantile estimé maximal : les $65$ paires emboîtées respectent l'ordre sur chaque graine. Ce constat n'est pas une preuve sur les valeurs intermédiaires des canaux ni un ordre stochastique général. Il n'est pas non plus trajectoriel. Enfin, rien ne garantit qu'une entité puisse atteindre simultanément les quatre valeurs hautes.}
\appui{supplément S1, \S~S1.4.}
\end{question}

\begin{question}{\og Pourquoi le pilier des tiers échappe-t-il à votre limite d'identification ?\fg{}}
\rep{Il n'échappe pas à cette limite. Le choc commun produit des co-occurrences sans direction, et les incidents chez les prestataires motivent ce mécanisme. Ils ne calibrent pas son intensité au sein d'une entité, ni la partie symétrique $S$, qui reste posée. La concentration des fournisseurs sert d'ancrage au scénario ; elle n'est pas une probabilité mesurée de transmission entre piliers.}
\appui{supplément S1, \S~S1.2, encadré.}
\end{question}

\subsection*{Supplément S4. Le protocole d'élicitation, diffusé et resté sans réponse}

\begin{parcoeur}
\begin{fichetab}
Cooke & quantiles $5$, $50$ et $95\,\%$ ; calibration sur des graines de réponse connue ;
  information ; poids égal au produit des deux, nul sous un seuil \\
Questionnaire & troisième version, seize questions dont dix graines (G1 à G10) et six liens (B1 à
  B6) ; valeurs vraies non publiées, pour ne pas brûler l'instrument ; graine $\xi=0{,}90$
  retirée (valeur posée, paramètre estimé) \\
Convention & $j$ vers $k$ ; une inversion de lettres ferait saisir $W^\top$, sans contrôle
  numérique possible \\
Diffusion & une commission de l'Institut et des praticiens du cabinet : aucune réponse, aucun
  panel \\
Pourquoi c'est tenable & à un ou deux experts, la pondération de Cooke dégénère ; un panel biaisé
  déplacerait le capital de $777$ M\euro{} sans signal, quand les bornes ne bougent pas \\
Relectures & trois réponses : zéro, deux et deux désaccords ; la cause commune avancée sans y être
  invité ; manquent les déclencheurs transversaux (gestion des changements, des
  vulnérabilités) \\
Coût & sans panel, P1 et P4 restent à égalité ($43{,}6$ contre $37{,}7\,\%$) ; reprise sans coût
  de développement avec une dizaine de répondants \\
\end{fichetab}
\end{parcoeur}

\begin{question}{\og Pourquoi n'avez-vous pas d'avis d'experts sur cette matrice ?\fg{}}
\rep{J'ai préparé le protocole, puis je l'ai envoyé à une commission de l'Institut et à des praticiens du cabinet. Je n'ai reçu aucune réponse, donc aucun panel n'a pu être constitué. Ce n'est pas un choix de ne pas essayer. Le mémoire chiffre aussi ce qu'un essai raté aurait coûté. Un panel biaisé déplace le capital de $777$~M\euro{} sans aucun signal d'alerte, alors que les bornes ne bougent pas. Une élicitation mal étalonnée échangerait une ignorance mesurée contre une erreur invisible.}
\appui{supplément S4.}
\end{question}

\begin{question}{\og Vous refusez le jugement d'expert, mais vous codez vous-même vos rapports. Contradiction ?\fg{}}
\rep{Je l'assume, pour trois raisons. Le codage est contrôlable : le protocole est écrit, les récits sont publics, et n'importe qui peut recoder. Son principal biais est mesuré par un test contre une loi nulle exacte. Et il donne un ordre partiel, que je borne ensuite. Une élicitation, elle, donnerait des coefficients qui entreraient directement dans le capital. Il reste une limite déclarée : un seul codeur.}
\appui{supplément S4.}
\end{question}

\begin{question}{\og Que vous ont dit les praticiens ?\fg{}}
\rep{Le premier ne contredit rien, et son exemple se lit de deux façons opposées, ce qui illustre la frontière. Le deuxième conteste deux phrases et propose de lui-même la cause commune. C'est la meilleure confirmation de la frontière. Le troisième répète les deux désaccords. Aucun ne distingue un ordre d'une cause commune. Je déclare ces retours sans les traiter, car les prendre en compte demanderait de recalibrer.}
\appui{supplément S4.}
\end{question}

\subsection*{Annexe B. Paramètres et notations}

\begin{parcoeur}
\begin{fichetab}
Trois natures & calibré, posé, semi-ancré ; la direction $u_{ij}$ est non identifiée \\
Forte sensibilité & seulement $\xi$, calibré, et les leviers voulus : $g$, le score de
  conformité, le multiplicateur de fréquence \\
Posés non leviers & copule ($\Delta$ de moins de $1{,}4\,\%$ sur $\theta\in[1,4]$), $\gamma$, seuil
  de crise $-2{,}5$, probabilités d'état, matérialité : influence faible à modérée \\
Notations & $g$ mesure une criticité, jamais un montant ; $\mathrm{VaR}$ et $\mathrm{TVaR}$ pour
  la charge annuelle $L$, $q$ et $\bar q$ pour un sinistre $X$ \\
\end{fichetab}
\end{parcoeur}

\begin{question}{\og Qu'est-ce qui est calibré, qu'est-ce qui est posé ?\fg{}}
\rep{Calibrés : la queue de sévérité, l'échelle, la fréquence, l'élasticité à la taille. Posés : la matrice et le gain de propagation, la surdispersion du moteur et la charge systémique (tous deux prudents), le plafond, les multiplicateurs de matérialité. Semi-ancrés : le multiplicateur de fréquence, la charge $\gamma$, les probabilités d'état. Aucune conclusion ne repose sur un paramètre à la fois non calibré et décisif hors de son rôle attendu.}
\appui{annexe~B, folio 91.}
\end{question}

\subsection*{Supplément S2. Pistes pour de futurs mémoires}

\begin{parcoeur}
\begin{fichetab}
Exécutables & exécuter l'élicitation préparée ; documenter les dépendances de tête, dont
  gouvernance-incidents ($11{,}2$ à $13{,}3\,\%$ chacune, ordre non résolu) ; le second codage en aveugle ; modéliser la dérive d'échelle mesurée \\
Donnée d'entité & calibrer l'amplitude des liens, pas seulement leur sens ; passer de l'échelle du
  secteur à une entité réelle \\
Structure & remplacer l'exposant d'additivité par un modèle de coût à contrainte de capacité ;
  étendre la cascade entre entités par le prestataire partagé, dont le graphe est observé \\
\end{fichetab}
\end{parcoeur}

\subsection*{Annexe C. Usage de l'IA générative}

\begin{parcoeur}
\begin{fichetab}
Usage & rédaction et reformulation du manuscrit ; écriture et relecture des scripts ; quatre
  agents de relecture automatique ; le dispositif de vérification lui-même \\
Vérification & chaque nombre publié est rattaché au script qui le produit, et un programme vérifie
  qu'il se retrouve dans les sorties versionnées de ce script ; il établit l'origine d'un nombre, non la
  justesse de la phrase \\
Limite & huit erreurs sémantiques que le dispositif confirmait toutes (cinq de relation, un
  renvoi, un superlatif, un comptage) ; deux venaient d'un message écrit en dur dans un script ;
  règle : un superlatif ou un comptage se calcule \\
Non délégué & la problématique, le choix de la cascade, la décision de borner, les sources et
  leurs biais, la hiérarchie des limites, l'interprétation \\
\end{fichetab}
\end{parcoeur}

\begin{question}{\og Avez-vous utilisé l'intelligence artificielle générative ?\fg{}}
\rep{Oui. C'est déclaré dans l'introduction et détaillé à l'annexe~C du rapport. J'ai utilisé l'IA pour m'aider à rédiger, pour écrire et relire les scripts, avec des agents de relecture et un dispositif de vérification. Aucun résultat chiffré n'en vient : chaque nombre publié sort de l'exécution d'un script du dépôt. Un dispositif rattache chaque nombre à la sortie qui le produit. L'annexe dit aussi ce que ce dispositif ne détecte pas.}
\appui{introduction, folio 11 ; annexe~C, folio 95.}
\end{question}

\begin{question}{\og Comment garantissez-vous qu'aucun chiffre n'est faux ?\fg{}}
\rep{Je ne le garantis pas, et l'annexe~C dit pourquoi. Les nombres vérifiables sont comparés aux sorties des scripts cités, et tous sont confirmés. Mais j'ai relu le texte en cherchant les superlatifs et les comptages, et j'ai trouvé huit erreurs de sens. Le dispositif les confirmait toutes : le nombre sortait du bon script, seule la phrase était fausse. Un taux de confirmation ne mesure pas la justesse du texte.}
\appui{annexe~C, folio 95.}
\end{question}

\begin{question}{\og Qu'est-ce que l'IA n'a pas fait ?\fg{}}
\rep{Tout ce qui engage. La problématique. Le choix de la cascade plutôt que d'une copule. La décision de ne pas poser de matrice arbitraire et de publier des bornes. Le traitement de l'élicitation restée sans réponse. Le jugement sur les sources. La hiérarchie des limites. L'interprétation. Et deux décisions vont à l'inverse d'une rédaction assistée : j'ai réfuté l'hypothèse qui donnait son titre au mémoire, et j'ai requalifié le chiffre d'entité en extrapolation dont la portée doit être déclarée.}
\appui{annexe~C, folio 95.}
\end{question}

% ---------------------------------------------------------------------
\section*{Questions transversales : la forme, les choix, la posture}
\addcontentsline{toc}{section}{Questions transversales : la forme, les choix, la posture}

\noindent Ces questions ne visent pas un chapitre mais les \emph{décisions} qui ont fait le mémoire,
et c'est là qu'un jury mesure la maturité. Chacune a une réponse en quelques phrases, et aucune ne
consiste à s'excuser.

\begin{question}{\og Votre rapport fait $100$~pages, les consignes parlent d'environ $30$~pages de corps.\fg{}}
\rep{Les consignes disent « environ », et la limite de $100$~pages, annexes et pièces liminaires comprises, vient du double usage du texte avec l'Institut des actuaires. Le volume vient des données, de la robustesse, et du fait que chaque nombre est rattaché à son script. J'ai déjà sorti ce qui pouvait l'être : les diagnostics secondaires, les adaptations par pilier et le questionnaire sont dans le supplément en ligne, sous DOI.}
\appui{notice des annexes, folio 86.}
\end{question}

\begin{question}{\og Pourquoi un supplément hors du rapport ?\fg{}}
\rep{Pour tenir les $100$~pages sans perdre la traçabilité. Les démonstrations propres à ma contribution, celles de la cascade, restent dans le rapport, en annexe~A. Sont sortis : les diagnostics secondaires, les adaptations par pilier, le questionnaire d'élicitation et les pistes. Le tout est déposé de façon durable sous l'identifiant \texttt{10.5281/zenodo.22874524}, donc on peut le citer. La version~1.3 est annoncée dans la notice des annexes.}
\appui{notice des annexes, folio 86 ; annexe~A, folio 87.}
\end{question}

\begin{question}{\og Vous citez une source que personne ne peut recalculer. Est-ce sérieux ?\fg{}}
\rep{Cette source donne la structure du risque, c'est-à-dire les parts par vecteur d'attaque. Elle ne donne jamais le niveau. Son statut est déclaré dans un encadré : je l'ai consultée, puis je ne l'ai pas conservée. Aucun script du dépôt ne peut donc la reproduire. J'ai préféré assumer une citation datée plutôt que fabriquer une reproductibilité que je n'ai pas. La retirer supprimerait une décomposition, pas une incertitude.}
\appui{\S~3.6, folio 27, encadré.}
\end{question}

\begin{question}{\og Vous exposez beaucoup vos limites. Ce n'est pas vous desservir ?\fg{}}
\rep{Je précise ce que chaque contrôle permet d'affirmer. Les calculs reproduisent un scénario sous des hypothèses déclarées ; cela ne valide pas ces hypothèses sur une entité réelle. Le chiffre d'entité reste une extrapolation, sans garantie de majoration. Une limite utile est celle dont je peux expliquer la conséquence et la donnée qui permettrait de la réduire.}
\appui{\S~10.2, folios 72 et 73 ; conclusion, folios 78 et 79.}
\end{question}

\begin{question}{\og Votre écart de $14\,139$~M\euro{} est-il connu au million près ?\fg{}}
\rep{Non, et le rapport ne le dit pas. L'écart est reproductible. Un retest indépendant sur seize nouvelles graines donne un facteur de $3{,}35$ contre $3{,}34$. Une inversion de Fourier donne $5\,945$ et $19\,890$, contre $6\,049$ et $20\,188$ en Monte-Carlo. Cet écart n'est pas significatif devant l'erreur type de la référence. Mais ce contrôle ne porte que sur le calcul du modèle une fois fixé. Les $734$~M\euro{} sont un écart-type entre quatre graines. L'erreur type de leur moyenne vaut $367$. L'incertitude sur la queue est bien plus large. En rééchantillonnant la donnée à seuil fixé (128 tirages, pertes supposées indépendantes), le rapport entre états va d'environ $3{,}1$ à $4{,}8$. C'est une bande exploratoire et conditionnelle, ni un intervalle de confiance ni des bornes générales. Je présente donc l'écart comme un résultat de scénario. Son ordre de grandeur est stable, et sa précision est limitée par la donnée.}
\appui{chapitre~9 pour les $734$ et $367$ ; \S~10.2 pour les limites ; retest du 20 octobre, hors rapport.}
\end{question}

\begin{question}{\og Pourquoi une perte indépendante par pilier touché, et non une par incident ?\fg{}}
\rep{C'est une hypothèse de construction, déclarée. Aucune source n'observe le coût d'un sinistre qui touche plusieurs piliers. Les pertes OpRisk donnent un montant et une année, sans ventilation par pilier. Elle pèse sur le résultat. Avec un seul tirage de perte par incident, le facteur entre états tomberait vers $2{,}2$ dans le retest indépendant. Mais ce contre-scénario supprime, par construction, l'effet du nombre de piliers sur le coût, donc le rôle financier de la propagation. La donnée n'identifie ni l'un ni l'autre. Le signe et l'ordre de grandeur de l'écart tiennent. Son amplitude est un scénario.}
\appui{\S~10.2, hypothèse d'additivité des coûts.}
\end{question}

\begin{question}{\og Que plafonne exactement le plafond de sévérité ?\fg{}}
\rep{Il plafonne chaque sévérité élémentaire, c'est-à-dire la perte d'un pilier touché. Il ne plafonne pas la perte totale de l'incident. Le code applique le plafond à chaque tirage, avant l'addition. Le plafond de $40$~M\euro{} de la source PRC et les diagnostics de plafond lié à l'exposition se lisent ainsi. Le résultat principal est calculé sur OpRisk et n'utilise pas de plafond.}
\appui{annexe~B, ligne du plafond ; \S~10.2.}
\end{question}

\begin{question}{\og Qu'auriez-vous fait avec six mois de plus ?\fg{}}
\rep{Trois choses, écrites au supplément~S2 avec ce qui existe déjà pour chacune. Je lancerais l'élicitation, qui est prête et ne manque que de répondants. Je ferais coder le corpus en aveugle par un second lecteur. Et pour la structure, je remplacerais l'additivité des coûts par un modèle avec contrainte de capacité. J'étendrais ensuite la cascade entre entités, via le prestataire partagé. Là, le graphe est \emph{observé}, donc l'extension n'hériterait pas du problème d'identification.}
\appui{supplément S2.}
\end{question}

\begin{question}{\og En une phrase, que faut-il retenir ?\fg{}}
\rep{La non-conformité à DORA a un coût en capital que la Formule Standard ne voit pas : un facteur $3{,}34$ au secteur, sur les quatre canaux. Et ce qu'on ne peut pas savoir, la direction de la contagion, on le borne et on le chiffre au lieu de le poser.}
\appui{conclusion, folios 78 et 79.}
\end{question}

\laphrase{Je n'ai pas cherché à rendre le chiffre plus sûr qu'il n'est. J'ai cherché à rendre exact
ce que l'on peut en dire, et à chiffrer ce que l'on ne peut pas.}

\begin{cle}
\textbf{La règle qui vaut pour toutes ces réponses.} Aucune ne consiste à défendre un chiffre :
toutes disent ce que le chiffre vaut et ce qu'il ne vaut pas. C'est la posture que le \S~4.6 du
référentiel valorise explicitement, et c'est celle qui a produit ce mémoire.
\end{cle}

% =====================================================================
\newpage
\section*{Annexe~: le modèle en formules, du sinistre au capital}
\addcontentsline{toc}{section}{Annexe : le modèle en formules}

\begin{cle}
\textbf{À quoi sert cette page.} Si l'on demande d'écrire le modèle au tableau, c'est cela qu'il
faut savoir poser, dans cet ordre, sans notes. Chaque bloc porte son \textbf{statut}, et ce statut
est la moitié de la réponse~: \emph{estimé} sur données, \emph{posé} faute de source,
\emph{démontré}, ou \emph{hypothèse de construction non testée}. Confondre les quatre est la
seule faute qui ne se rattrape pas.
\end{cle}

\subsection*{1. La chaîne, en une ligne}

\[
\text{fréquence} \;\longrightarrow\; \text{cascade entre piliers} \;\longrightarrow\;
\text{coût du sinistre} \;\longrightarrow\; \text{charge annuelle} \;\longrightarrow\;
\text{capital}
\]

\noindent L'état de conformité n'ajoute aucune pénalité~: il \textbf{déplace quatre paramètres} de
cette chaîne, et rien d'autre.

\subsection*{2. La fréquence \quad\small\textcolor{navy}{statut~: estimée}\normalsize}

Le nombre annuel d'incidents enregistrés suit une binomiale négative, retenue contre la loi de
Poisson parce que la variance observée excède la moyenne~:
\[
N \sim \mathcal{BN}(r,p), \qquad \mathbb{E}[N]=\lambda, \qquad
\frac{\mathrm{Var}(N)}{\mathbb{E}[N]} \;=\; 1+\frac{\lambda}{r}.
\]
\lepiege{L'indice de dispersion \textbf{n'est pas invariant d'échelle}~: il dépend de $\lambda$,
donc il vaut autre chose sur des cellules firme-année, sur la série agrégée et au moteur. Trois
valeurs du projet sont trois \emph{échelles de lecture}, pas trois mesures contradictoires.}

\subsection*{3. La sévérité \quad\small\textcolor{navy}{statut~: estimée}\normalsize}

Au-delà du seuil $u$, les excès suivent une loi de Pareto généralisée, ce que le théorème de
Balkema et de Haan justifie asymptotiquement~:
\[
\mathbb{P}(X>u+y \mid X>u) \;=\; \Bigl(1+\xi\,\frac{y}{\sigma}\Bigr)^{-1/\xi},
\qquad y\ge 0 .
\]
D'où le quantile de \emph{sévérité d'un sinistre}, avec $\zeta_u=\mathbb{P}(X>u)$~:
\[
q_\alpha \;=\; u \;+\; \frac{\sigma}{\xi}
\left[\Bigl(\frac{1-\alpha}{\zeta_u}\Bigr)^{-\xi}-1\right].
\]
\cheznous{$u=20{,}03$~M\euro, $91$ excès, $\xi=0{,}5954$, $\sigma=57{,}97$.}
\lepiege{$\xi$ n'est pas qu'un paramètre de forme, c'est une \textbf{condition d'existence}~: la
moyenne existe si $\xi<1$, la variance si $\xi<1/2$. C'est cet argument, et non un argument de
degré, qui écarte deux variantes de modèle.}

\subsection*{4. La cascade \quad\small\textcolor{navy}{statut~: normalisation démontrée, amplitude posée}\normalsize}

On part d'une matrice d'influence brute $T$, asymétrique, et on la normalise à la
Leontief par l'\textbf{émission} maximale~:
\[
s_j=\sum_k T_{jk}, \qquad s_{\max}=\max_j s_j, \qquad
W \;=\; g\,\frac{T}{s_{\max}}, \qquad
e_j \;=\; g\,\frac{s_j}{s_{\max}} \;\le\; g .
\]
\noindent La convention est que $W_{jk}$ va de la \textbf{source} $j$ vers la \textbf{cible} $k$~:
la ligne émet, la colonne reçoit. La normalisation \emph{garantit} que le pilier le plus prolifique
engendre au plus $g$ descendants directs, donc $\rho(W)\le g<1$ et la sous-criticité est acquise
par construction, non supposée.

\cheznous{$s_{\max}=2{,}60$, $\rho(W)=0{,}506$, $R_0=0{,}054$, et le $g$ critique vaut $1{,}78$,
hors du domaine admissible.}
\lepiege{Diviser par la \emph{réception} maximale au lieu de l'émission donnait $e_{P_1}=1{,}02$
pour $g=0{,}9$~: la borne que la normalisation existe pour poser était franchie, et $g$ cessait de
désigner ce qu'il désigne.}

\subsection*{5. L'ensemble touché, et le coût d'un sinistre}

Le moteur tire une \textbf{marche auto-évitante}~: depuis le pilier courant $j$, la cascade
continue avec probabilité $e_j$, le successeur est tiré proportionnellement à
$T_{j\cdot}$, et retomber sur un pilier déjà touché l'éteint. On obtient l'ensemble
$S_a$ des piliers atteints depuis l'amorce $a$, et le coût du sinistre vaut
\[
C \;=\; \sum_{j\in S_a} c_j
\qquad\quad
\small\textcolor{accent}{\textbf{hypothèse de construction, non testée}}\normalsize
\]
\lepiege{C'est la plus lourde des hypothèses du mémoire. Relâchée par un exposant sur $|S_a|$, elle
déplace l'écart de capital de $9\,706$ à $20\,487$~M\euro. Ce qui survit à toute la plage~: le
signe et l'ordre. Ce qui ne se déclare pas~: le \emph{sens} de l'effet.}

\subsection*{6. La charge annuelle et le capital \quad\small\textcolor{navy}{statut~: définition}\normalsize}

\[
L \;=\; \sum_{i=1}^{N} C_i, \qquad\qquad
\mathrm{SCR} \;=\; \mathrm{VaR}_{99{,}5\,\%}(L)\quad\text{(non centrée)}.
\]
\lepiege{Le capital publié est le quantile \emph{entier}, sans retrancher $\mathbb E[L]$ : la
perte attendue d'une non-conformité n'est couverte par aucune provision, donc elle n'en est pas
déduite. Centré, le chiffre de tête tomberait à environ $5\,430\to16\,320$, un facteur de l'ordre
de $3{,}0$ (calcul hors mémoire).}
\lepiege{$\mathrm{VaR}$ et $\mathrm{TVaR}$ sont réservées à la charge \textbf{annuelle agrégée}
$L$~; $q_\alpha$ du bloc~3 est un quantile de \emph{sévérité d'un sinistre}. Les confondre fait
lire une différence d'échelle comme une contradiction, et c'est arrivé en séance. Le niveau
$99{,}5\,\%$ n'a de sens réglementaire que sur un horizon annuel~: une perte isolée n'en porte
aucun.}

\subsection*{7. Les états de conformité \quad\small\textcolor{navy}{statut~: structure posée, charge posée}\normalsize}

La conformité de chaque pilier dérive d'une latente à facteur commun, transposée du risque de
crédit~:
\[
C^\star_j \;=\; \gamma\,\Theta \;+\; \sqrt{1-\gamma^{2}}\;\varepsilon_j,
\qquad \Theta,\varepsilon_j \ \text{indépendantes} \ \mathcal{N}(0,1),
\]
\[
\mathrm{corr}(C^\star_j, C^\star_k) \;=\; \gamma^{2} \;=\; 0{,}4624
\quad\text{pour } \gamma=0{,}68 .
\]
\noindent Des seuils ordonnés $K_{\text{bas}}<K_{\text{haut}}$ découpent les états non conforme,
partiel et conforme.
\lepiege{La corrélation de $0{,}462$ \textbf{n'est pas posée}~: elle est le \emph{carré} de la
charge de facteur commun. Si l'on demande d'où elle sort, c'est la réponse.}

\subsection*{8. L'identification partielle \quad\small\textcolor{navy}{statut~: démontré}\normalsize}

Toute matrice se décompose de façon unique en une part symétrique et une part antisymétrique,
orthogonales au sens de Frobenius~:
\[
W = S + A, \qquad S=\tfrac12\,(W+W^{\!\top}), \qquad A=\tfrac12\,(W-W^{\!\top}).
\]
\noindent La donnée identifie $S$, la co-occurrence~; elle n'identifie pas $A$, la direction. On
ne pose donc pas $A$, on la \textbf{borne}, par un degré d'ignorance $t$~:
\[
\mathcal{W}(t) \;=\; \bigl\{\,W=S+A \ :\ |A_{jk}| \le t\,S_{jk}\,\bigr\},
\qquad
\bigl[\underline{\mathrm{SCR}},\ \overline{\mathrm{SCR}}\bigr]
= \bigl[\min_{\mathcal{W}(t)}, \ \max_{\mathcal{W}(t)}\bigr].
\]
\cheznous{Cinq piliers donnent dix paires, donc $2^{10}=1\,024$ sommets, énumérés
exhaustivement. Bornes $[6\,858\,;8\,697]$~M\euro, largeur $1\,839$.}
\laphrase{Le paramètre $t$ élargit les asymétries admises. Leur effet sur le capital est calculé aux sommets, conditionnellement à la force symétrique posée. Cette plage ne mesure pas une incertitude d'échantillonnage et n'est pas démontrée linéaire en $t$.}

\subsection*{9. Les quatre canaux et leur interaction \quad\small\textcolor{navy}{statut~: mesuré}\normalsize}

La non-conformité déplace quatre paramètres, et l'on énumère les $2^4=16$ configurations. Pour un
sous-ensemble de canaux $S$, l'effet propre se lit par une décomposition de Möbius~:
\[
m(S) \;=\; \sum_{T\subseteq S} (-1)^{|S|-|T|}\,\Delta(T),
\qquad
\Delta(\text{tous}) \;=\; \sum_{S} m(S) .
\]
\cheznous{$\lambda$ de $21{,}6$ à $53{,}6$, $\zeta_u$ de $0{,}128$ à $0{,}181$, $g$ de $0{,}45$ à
$0{,}90$, accumulation tiers de $0$ à $0{,}68$. Capital de $6\,049$ à $20\,188$~M\euro, facteur
$3{,}34$, écart $14\,139$. Les canaux isolés somment à $9\,138$~: il manque $5\,001$~M\euro{}
d'interaction, soit $+35\,\%$.}
\lepiege{\textbf{Ne jamais additionner les quatre effets isolés} pour chiffrer une remédiation
partielle. Et ne pas confondre cette interaction, entre \emph{canaux}, avec celle entre
\emph{piliers}~: deux partitions du même écart, et la seconde a un signe non résolu.}

\subsection*{10. Le pont entre les deux échelles}

Question classique~: pourquoi un quantile de sévérité de quelques centaines de millions
cohabite-t-il avec un capital de plusieurs milliers de millions~? Parce qu'à $\lambda\simeq21{,}6$
sinistres par an, atteindre le quantile \emph{annuel} à $99{,}5\,\%$ demande un quantile \emph{par
sinistre} de l'ordre de $99{,}98\,\%$, bien plus haut dans la queue. L'approximation par perte
unique en donne la forme, la survie de la somme valant asymptotiquement $\mathbb E[K]$ fois celle
d'une perte isolée :
\[
\mathrm{VaR}_\alpha(L)\;\approx\;q_{\,1-(1-\alpha)/\mathbb E[K]}(X),
\qquad\text{d'où un quantile qui croît comme } (\lambda\,p_u)^{\xi}.
\]
\cheznous{Pour le canal de fréquence, cette forme prédit un facteur
$(53{,}62/21{,}56)^{0{,}5954}=1{,}720$ entre états, contre $1{,}715$ simulé ; pour la détection,
$1{,}228$ contre $1{,}239$. Au second ordre, elle retrouve le quantile publié à moins de $10\,\%$.}
\lepiege{Le quantile de sévérité ($663$~M\euro) n'est pas un capital, et le rapport entre les deux
objets n'est pas fixe : à l'échelle d'entité, le capital ($169$~M\euro) passe même sous le quantile
unitaire. Les comparer ne conclut jamais, dans un sens comme dans l'autre.}

% =====================================================================
\newpage
\section*{Annexe~: les nombres à connaître par c\oe ur}
\addcontentsline{toc}{section}{Annexe : les nombres à connaître}

% TABLE RESSERREE LE 23 SEPTEMBRE 2026. A l'interligne de 1,25 elle debordait sa page de
% 67,8 pt, un tabular ne sachant pas se couper : le bas de la table sortait du cadre sans
% que rien ne le signale a la lecture du source. Defaut anterieur, trouve a la compilation.
% PASSEE EN \footnotesize LE 30 SEPTEMBRE 2026 : en \small, avec les lignes ajoutees, l'encadre
% final tombait seul sur une page.
\begingroup\footnotesize
\renewcommand{\arraystretch}{1.02}
\begin{longtable}{@{}p{5.0cm}p{9.2cm}p{1.1cm}@{}}
\toprule
\textbf{Nombre} & \textbf{Ce que c'est} & \textbf{Ch.} \\
\midrule
\multicolumn{3}{@{}l}{\textcolor{navy}{\textbf{Le chiffre de tête, à savoir sans hésiter}}} \\
$6\,049 \to 20\,188$~M\euro & besoin de capital, conforme vers non conforme, au secteur & 9 \\
$3{,}34$ & le facteur entre les deux états & 9 \\
$14\,139$~M\euro & l'écart, dont $9\,138$ portés par les canaux isolés & 9 \\
$5\,001$~M\euro\ ($35\,\%$) & la part portée par l'interaction des canaux & 9 \\
$10\,377$ / $11\,413$ & la fréquence seule / les trois autres canaux réunis & 9 \\
$[6\,858\,;\,8\,697]$~M\euro & \textbf{la réponse à \og quel est votre SCR\fg{}} & 8 \\
$1\,839$~M\euro & largeur de plage aux sommets, à $S$ fixée & 8 \\
\midrule
$582$ / $91$ & pertes cyber $\times$ finance / excès au-delà du seuil & 3, 4 \\
$u=20{,}03$~M\euro & le seuil des excès (percentile $84{,}4$) & 4 \\
$\xi=0{,}5954$ / $1{,}03$ & indice de queue OpRisk / PRC & 4 \\
$663$ / $1\,752$~M\euro & quantile et moyenne de queue d'\emph{un} sinistre, pas un capital & 4 \\
$[411\,;1\,037]$, $\times2{,}5$ & intervalle bootstrap de ce quantile & 4 \\
$p_u=0{,}1509$ & $86/570$ (2000--2026), contre $91/582$ (complet) & 10 \\
$341$ / $21{,}6$ / $0{,}092$ & les trois fréquences (facteur $3\,719$) & 3 \\
$1{,}19$ / $2{,}24$ / $9{,}20$ & surdispersion : cellules, série, moteur (posé) & 4 \\
$p\approx10^{-30}$ & rejet de Poisson contre la binomiale négative & 4 \\
$0{,}087$ $[0{,}026\,;0{,}148]$ & élasticité sévérité-taille, corrigée du seuil & 3 \\
\midrule
$g=0{,}45$ / $0{,}68$ / $0{,}90$ & gains de propagation posés, conforme / partiel / non conforme & 6 \\
$\rho(T)=1{,}461$ & pourquoi il faut renormaliser & 5 \\
$\rho(W)=0{,}506$, $R_0=0{,}054$ & les deux rayons spectraux & 5 \\
$62{,}31\,\%$ / $31{,}15\,\%$ & multi-piliers NC/C hors P4 ; $73{,}25\,\%$ NC avec P4 & 10 \\
$R^2=0{,}9945$ & le capital, presque additif en piliers & 6 \\
$K_4=1{,}03$ ; $\lambda_U=0{,}53$ & seuil de P4 ; dépendance de queue de Gumbel & S1 \\
\midrule
$119$ contre $128\pm27$ & placebo sur OpRisk, $z=-0{,}33$ & 7 \\
$z=+3{,}93\to+5{,}12$ & placebo sur $7$ puis $10$ rapports & 7 \\
$0{,}12\,\%$ / $1{,}19\,\%$ / $1$ & VERIS : champ, chronologie, intersection & 7 \\
$12$ / $0$ sur $20$ & liens $p_{jk}$ crédiblement faibles / forts & 7 \\
$5\,275$~M\euro & socle sans contagion & 8 \\
$777$~M\euro & ce qu'un panel biaisé déplacerait sans signal & 8 \\
\midrule
$8\,123$ / $169$~M\euro & besoin sectoriel / extrapolation d'entité & 9 \\
$47$ contre $8{,}3$ & multiple de capital, entité contre secteur & 9 \\
$2{,}0$ à $26{,}4\,\%$ & charge rapportée au SCR publié, quatre SFCR & 9 \\
$-64\,\%$ ; $-56\,\%$ ; $64\,\%$ & gain du traité hors / pertes conservées comprises ; seuil de portée sans allègement & 8 \\
$4{,}75\,\%$ & coût du capital, directive (UE) 2025/2 & 9 \\
$657$ / $645$ / $1\,037$ & point, prédictive, robuste à $95\,\%$ & 10 \\
$\times8{,}93$ / $\times1{,}27$ / $\times9{,}21$ & paramètre, identification, famille de queue & 10 \\
$91{,}7\,\%$ / $75{,}0\,\%$ & backtest : binomiale négative / Poisson & 4 \\
$4{,}00\,\%$ par an & dérive d'échelle de la queue & 4 \\
$1\,811$ années & pour tester $99{,}5\,\%$ à $80\,\%$ de puissance & 4 \\
$32{,}7\,\%$ / $13{,}3\,\%$ & valeur d'information moyenne : six sens / paire de tête & 10.8 \\
\midrule
\multicolumn{3}{@{}l}{\textcolor{navy}{\textbf{Robustesse du classement, à sortir si l'on attaque le corpus}}} \\
$43{,}6\,\%$ / $37{,}7\,\%$ & P1 / P4 premiers sur l'ensemble admissible & 9 \\
$1{,}5\,\%$ & écart de regret maximal entre les deux & 9 \\
$a=0{,}25$ / $a=0{,}50$ & le classement tient / bascule vers P4 & 7, 9 \\
$-24{,}8\,\%$ ($0{,}81\times$) & chute du niveau d'entité si P1 devient muette & 7 \\
$34\,\%$ & part de Shapley de P1 dans le surcoût & 9 \\
\bottomrule
\end{longtable}
\endgroup

\vspace{4mm}
\begin{cle}
\textbf{La phrase qui tient tout.} Si l'on ne devait retenir qu'une chose de ce document~:
\emph{le niveau de capital bouge dans une bande large ; ce qui résiste, c'est le signe et l'ordre
de grandeur de l'écart entre états face à un forfait qui ne bouge pas, le rapport au socle, et le
groupe le plus souvent en tête dans les directions examinées, P1 et P4, sans garantie universelle}. L'ordre entre P1 et P4, lui, dépend
d'une direction que la donnée ne fixe pas, et le dire fait partie de la réponse.
\end{cle}


% Réponses complémentaires actualisées le 2 novembre 2026 ; sorties 136--139.
% Ne pas attribuer ces calculs au rapport ENSAE déposé. Nom conservé pour compatibilité.
\clearpage
\part*{Compléments récents : questions Institut après dépôt}
\addcontentsline{toc}{part}{Compléments récents : questions Institut après dépôt}
\begin{cle}
Les questions suivantes portent sur les calculs Institut disponibles au 2 novembre. Pour l'ENSAE,
les présenter seulement comme des compléments après dépôt. Les appuis sont des sections Institut,
des scripts et leurs sorties ; les folios ENSAE du reste du document sont conservés.
\end{cle}

\begin{question}{\og Votre résultat est-il 169 ou 24 millions, et est-ce une borne prudente ?\fg{}}
\rep{Le rapport ENSAE déposé conserve $169$~M\euro{}, lecture historique du registre. La variante
retenue dans l'Institut donne $24{,}1$~M\euro{} pour la notionnelle de $20$~Md\$ d'actifs, dans
l'état supposé non conforme. Fréquence, dispersion et corps sous le seuil changent ensemble.
Le chiffre n'est pas une borne prudente ni un capital individuel validé. La plage aux sommets
est $15{,}8$ à $27{,}5$~M\euro{} ; le bootstrap des firmes donne $8{,}3$ à $53{,}0$, sans corriger
la sélection du registre. Les valeurs aux sommets ne sont pas des extrema globaux démontrés.}
\appui{Institut, \S~9.5--9.6 ; scripts 136 et 137.}
\end{question}

\begin{question}{\og Pourquoi retenir 0,494 plutôt que 0,074 ou 0,27 ?\fg{}}
\rep{Le $0{,}074$ est ajusté sur tout le registre, où les petites firmes ont des fenêtres
d'observation courtes. Le $0{,}27$ était un milieu géométrique posé entre ce coefficient et
la proportionnalité. La variante retenue ajuste une NB2 sur $1\,402$ firmes de plus de
$100$~Md\$ : pente $0{,}494$, erreur-type par firme $0{,}066$, forme $r=0{,}594$.
L'ajustement donne $\lambda(A)=0{,}0209$/an à $20$~Md\$ d'actifs. La restriction ne prouve
pas que le biais de sélection est corrigé : $57{,}8\,\%$ des grandes firmes ne sont suivies
qu'un an. L'extrapolation en dessous de $100$~Md\$ reste la limite principale.}
\appui{Institut, équation~9.1 et \S~9.5 ; scripts 136 et 137.}
\end{question}

\begin{question}{\og Comment votre VaR peut-elle être nulle alors que des incidents surviennent ?\fg{}}
\rep{Le quantile porte sur la perte annuelle. Sous la lecture retenue, le corps empirique est
strictement positif : un incident produit une perte, mais sa probabilité annuelle peut rester
inférieure à $0{,}5\,\%$. La petite entité de $500$~M\euro{} d'actifs a une VaR nulle et une
ES d'environ $23$~M\euro{}. L'assureur A a désormais une VaR d'environ $1$~M\euro{}.
La bascule en forme fermée est $1\,076$~M\euro{} d'actifs, avec une plage bootstrap de
$100$ à $3\,889$~M\euro{}. Dans l'ancien moteur, une perte pouvait en plus être nulle sous le seuil ;
ce mécanisme ne décrit pas la variante actuelle. Une VaR nulle ne signifie pas risque nul.}
\appui{Institut, \S~9.6 ; scripts 136 et 137.}
\end{question}

\begin{question}{\og Un million pour l'assureur A, n'est-ce pas trop faible ?\fg{}}
\rep{Je ne peux pas valider ce niveau. La probabilité de perte n'est que de $0{,}74\,\%$ dans
le scénario de référence : la VaR à $99{,}5\,\%$ est proche de la masse en zéro. L'ES vaut
environ $51$~M\euro{}. Si la capture supposée passe de $100$ à $50\,\%$, la VaR de A passe
de $0{,}9$ à $8{,}9$~M\euro{} sur quatre graines. Ce diagnostic rend visible la sensibilité
à l'observation ; il ne permet ni de choisir un taux de capture ni d'affirmer que le besoin
réel est couvert.}
\appui{Institut, tableau de comparaison et tableau 9.6 ; scripts 136 et 138.}
\end{question}

\begin{question}{\og Comment savez-vous combien d'incidents ne sont pas enregistrés ?\fg{}}
\rep{Je ne le sais pas. Une capture binomiale de probabilité $q_{\mathrm{obs}}$ transforme
une NB de moyenne $\mu/q_{\mathrm{obs}}$ en une NB de moyenne observée $\mu$, sans changer
sa forme $r$. Plusieurs risques réels ont donc la même loi observée. Pour la notionnelle,
les captures posées de $100$, $50$, $25$ et $10\,\%$ donnent $23{,}1$, $74{,}0$, $153{,}9$
et $327{,}7$~M\euro{}. Ce sont des sensibilités, pas des estimations. La capture est supposée
indépendante du montant ; si les grosses pertes sont plus publiées, sévérité et capture doivent
être modélisées ensemble. Le périmètre financier souffre de cette limite lui aussi.}
\appui{Institut, \S~9.5.2 et 9.6.2 ; script 138.}
\end{question}

\begin{question}{\og Que fait exactement le passage conforme / non conforme dans les calculs d'entité ?\fg{}}
\rep{Dans ce diagnostic d'entité, seul le gain de propagation change : $g=0{,}45$, $0{,}68$
et $0{,}90$ donnent $13{,}2$, $17{,}7$ et $24{,}1$~M\euro{}, soit $+83\,\%$ entre C et NC.
Le chiffre à quatre canaux du périmètre financier fait aussi varier fréquence, détection et
accumulation. Ce ne sont pas les mêmes protocoles, et le facteur $3{,}34$ ne se transpose pas
directement à l'entité. La figure 9.3 utilise désormais le protocole d'entité du script 136,
contrôlé et retracé par le script 139.}
\appui{Institut, \S~9.1 et 9.5, figure 9.3 ; scripts 136 et 139.}
\end{question}

\begin{question}{\og Quels paramètres sont estimés, posés ou bornés ?\fg{}}
\rep{La fréquence de base est comptée dans le registre et la loi de sévérité est ajustée
aux pertes. Pour la fréquence non conforme, les parts des vecteurs d'attaque sont observées ;
les multiplicateurs sont les centres de fourchettes traduites de résultats publiés sur des
contrôles. Ce passage est semi-ancré, non estimé directement. Les facteurs de détection,
les gains par état, la structure relative $T$ et la probabilité de choc commun sont posés.
Les fourchettes et les contraintes sur les directions de $W$ définissent des scénarios
admissibles, pas des intervalles de confiance. La co-occurrence n'est pas estimée ici.}
\appui{Beamer ENSAE, diapo 9 ; Institut, chapitres 4 à 6 et annexe E.}
\end{question}

\begin{question}{\og Le placebo MOVEit prouve-t-il un effet net nul ?\fg{}}
\rep{Non. Les $+1{,}17$ jours et le placebo de $+1{,}11$ donnent une différence descriptive
de $+0{,}06$. Le premier calcul porte sur $188$ organisations traitées, le placebo sur $20$ ;
aucun intervalle ni test de leur contraste n'est calculé. Le placebo fragilise l'interprétation
causale du dispositif ; il ne démontre ni un effet net nul ni une absence générale de contagion.}
\appui{Institut, \S~7.4 ; script 35.}
\end{question}

\begin{question}{\og La gouvernance et les tiers sont-ils toujours les deux priorités ?\fg{}}
\rep{Ils sont prioritaires au point d'expert et arrivent le plus souvent en tête lorsque
la direction varie. Dans les $220$ directions examinées, P1 et P4 sont premiers dans
$43{,}6\,\%$ et $37{,}7\,\%$ des cas ; P2 ou P3 peuvent passer devant. Une contribution
de Shapley, le gain de remédier un pilier et le premier d'un ordre optimal sont des objets
différents. Je précise la matrice, l'état de départ et le critère de décision avant de recommander.}
\appui{Institut, \S~9.3 et chapitre 10 ; sortie 30 pour les gains de remédiation.}
\end{question}

\begin{question}{\og Tous vos résultats utilisent-ils exactement la même cascade ?\fg{}}
\rep{Non. Le chiffre de tête à quatre canaux emploie la marche auto-évitante avec un choc commun
P4. Les configurations par pilier utilisent la marche sans ce choc. L'exploration des matrices
et les calculs d'entité utilisent la cascade indépendante sans choc commun. Je précise le
protocole avant de comparer les montants ; une vérification sur l'un ne valide pas automatiquement les autres.}
\appui{Institut, tableau des mécanismes du chapitre 5 ; \S~6.5 et \S~9.5.}
\end{question}

\begin{question}{\og Pourquoi ne pas remplacer votre modèle par de l'apprentissage statistique ?\fg{}}
\rep{Je l'ai testé sur une question précise : les caractéristiques améliorent-elles la prévision
du dépassement de seuil parmi les pertes enregistrées ? Sur $280$ pertes évaluées en origine
glissante, le boosting a le meilleur score moyen. Son gain contre la tendance seule est
$+0{,}0193$, avec un intervalle à $90\,\%$ de $-0{,}0371$ à $+0{,}0571$. Le gain n'est donc
pas résolu. Cela ne prouve pas une impossibilité ; avec ce protocole et ces données, je n'ai
pas de justification empirique suffisante pour remplacer la référence.}
\appui{Institut, \S~4.7.1 ; script 131.}
\end{question}

\begin{question}{\og Vos actifs ne contiennent-ils pas une information postérieure à l'incident ?\fg{}}
\rep{Oui, pour $37\,\%$ des pertes. La version déclarée utilise l'année d'information financière
disponible dans la base, ce qui peut introduire une fuite temporelle. Le diagnostic masque
donc les actifs postérieurs à l'année de l'événement, et compare aussi des modèles sans actifs.
Les intervalles des gains incluent encore zéro. La validation repose toutefois sur douze blocs
annuels, et certaines firmes sont présentes avant et après la séparation temporelle.}
\appui{Script 131, disponibilités des actifs et sensibilités.}
\end{question}

\begin{question}{\og Votre GPD enrichie par les caractéristiques prédit-elle mieux les excès ?\fg{}}
\rep{Aucun gain résolu n'apparaît, et il faut citer l'échec de support : la variante enrichie
donne une densité nulle à une perte observée. La moyenne de score sur les $59$ autres excès
omet cette observation ; ce n'est donc pas son score complet. La comparaison est utile pour
tester une alternative, mais elle ne valide pas la GPD commune ni son quantile annuel.}
\appui{Institut, \S~4.7.1 ; partie sévérité du script 131.}
\end{question}

\begin{question}{\og Les descriptions SAS permettent-elles enfin de calibrer votre matrice ?\fg{}}
\rep{Pas à ce stade. Le codage complet des $570$ descriptions donne $65$ incidents multi-piliers
au sens strict, $203$ au sens large. L'accord de rang avec l'expert paraît élevé sur les
co-occurrences brutes, mais devient faible après prise en compte des marges : $0{,}33$ et
$0{,}12$. Un seul assistant IA a codé les textes, sans validation indépendante ; aucun sens
de transmission n'est codé. Ces résultats n'estiment ni $S$ ni $A$.}
\appui{Script 133 ; le script 132 est le pilote de 50 descriptions, inclus dans ce corpus.}
\end{question}

\begin{question}{\og Qu'avez-vous préparé pour vérifier ce codage ?\fg{}}
\rep{Un lot en aveugle de $100$ descriptions, équilibré entre $50$ fraudes et $50$ interruptions.
Le second lecteur voit les textes et la grille, sans les codes du premier lecteur. L'outil
refuse de traiter une réponse manquante comme un ensemble vide, puis calcule les accords et
le kappa par pilier et par strate. Le lot n'a pas les proportions de la base ; les réponses
indépendantes sont encore attendues. Aucun chiffre d'accord n'est donc présenté.}
\appui{\path{docs/protocoles/validation_codage_sas.md} ; outil \texttt{double\_codage\_sas.py}.}
\end{question}


\end{document}
